\chapter{Electromagnetism: Maxwell's equations as four observations}
\label{vol04:ch08}

\epigraph{The agreement of the results seems to show that light and
magnetism are affections of the same substance, and that light is an
electromagnetic disturbance propagated through the field according to
electromagnetic laws.}{James Clerk Maxwell, \emph{A Treatise on
Electricity and Magnetism}, 3rd ed.\ (Oxford: Clarendon, 1892),
vol.~II, art.~781 \cite[vol.~II, p.~431]{hist:maxwell1873}.}

\chapmeta{Half-life: durable. Archetypes: balance, transport, interface. Exercise target: 31. Project track: physical measurement plus analysis (hybrid). Hazard class: low (bench-scale currents and voltages, mains avoided). Reader path default: standard, with sections explicitly tagged. Named cases: lightning-induced surge damage to grid and electronic equipment (described by mechanism, no single accident named).}

\noindent
The four Maxwell equations are observations of the world before they
are mathematics. Each was discovered as a regularity in experiment, and
each retains the working content of its experimental origin. We state
the four equations in their integral form, derive their differential
form, show that they imply the wave equation, and end with the
failure modes of every electrical system: ground loops, induced
electromagnetic interference, and lightning. The reader who finishes
the chapter writes energy balances across electromagnetic fields with
the same discipline applied in earlier volumes to mass and momentum.

\section{Charge, current, and fields}\pathtag{core}

Electric charge is a property of matter measured in coulombs, $1\,\si{
\coulomb} = 1\,\si{\ampere\second}$. The charge of a proton is
$+e = 1.602\,176\,634\times 10^{-19}\,\si{\coulomb}$, fixed exactly in
the 2019 SI redefinition; the charge of an electron is $-e$. Macroscopic
objects carry net charge equal to the algebraic sum of their elementary
charges. Charge is conserved: the total charge inside any closed
surface changes only by the flux of charge through the surface.

\engterm{Current} is the rate at which charge crosses an oriented
surface,
\[
I = \dv{Q}{t}.
\]
For a wire carrying $1\,\si{\ampere}$, about $6.24\times 10^{18}$
elementary charges cross any cross-section per second. The drift speed
of conduction electrons in a typical copper wire is on the order of
$10^{-4}\,\si{\meter\per\second}$; the electrons are slow, but their
density is high.

The \engterm{current density} $\vect{J}$ is the local current per unit
area, with direction along the flow,
\[
I = \int_{S} \vect{J}\cdot \dd\vect{A},
\]
so $\vect{J}$ carries units $\si{\ampere\per\meter\squared}$. Charge
conservation, applied to an infinitesimal control volume, gives the
\engterm{continuity equation}:
\[
\pdv{\rho}{t} + \divg \vect{J} = 0,
\]
where $\rho$ is the charge density in $\si{\coulomb\per\meter\cubed}$.
The continuity equation is the same accounting equation that closed the
mass balance in Volume IV, Chapter 7: stuff in minus stuff out equals
change in stuff stored, with $\vect{J}$ playing the role of mass flux
and $\rho$ the role of mass density. Every Maxwell equation can be read
as a balance once the right field is identified as flux and the right
field as source.

A \engterm{field} is a quantity defined at every point of space. The
electric field $\vect{E}(\vect{r}, t)$ is the force per unit charge that
a test charge would experience at $\vect{r}$ at time $t$,
\[
\vect{F} = q\vect{E},
\]
in units of $\si{\volt\per\meter}$ (equivalently $\si{\newton\per
\coulomb}$). The magnetic field $\vect{B}(\vect{r}, t)$ is the
proportionality between velocity and the velocity-dependent component
of the force on a charge,
\[
\vect{F} = q\vect{v}\times \vect{B},
\]
in units of teslas, $\si{\tesla} = \si{\volt\second\per\meter\squared}$.
The Earth's magnetic field at the surface is around $30$ to $60\,\si{
\micro\tesla}$; a refrigerator magnet is around $5\,\si{\milli\tesla}$;
a clinical MRI machine is $1.5$ to $7\,\si{\tesla}$. Both fields are
observable through the forces they exert; both are real, in the sense
that they store energy that can do work.

% Figure: electric field lines of a single positive point charge and of an
% equal-and-opposite pair (a dipole). The left panel shows radial lines
% diverging from a positive charge; the right panel shows the dipole pattern
% running from the positive to the negative charge.
\begin{figure}[htbp]
\centering
\begin{tikzpicture}[>=Latex, scale=1.0]
% === Left panel: isolated positive point charge, radial field lines ===
\begin{scope}[shift={(0,0)}]
  \foreach \a in {0,30,...,330}{
    \draw[engblue, ->] (\a:0.35) -- (\a:1.7);
  }
  \filldraw[engred] (0,0) circle [radius=4pt];
  \node[white, font=\scriptsize\bfseries] at (0,0) {$+$};
  \node[font=\small] at (0,-2.2) {isolated positive charge};
\end{scope}
% === Right panel: dipole ===
\begin{scope}[shift={(6.2,0)}]
  \coordinate (P) at (-0.9,0);
  \coordinate (N) at (0.9,0);
  % a handful of representative field-line arcs from + to -, bowing by a
  % fixed offset so no trigonometric overflow can occur
  \foreach \h in {0.45,1.0,1.7}{
    \draw[engblue] (P) .. controls (-0.4,\h) and (0.4,\h) .. (N);
    \draw[engblue] (P) .. controls (-0.4,-\h) and (0.4,-\h) .. (N);
  }
  % arrowheads on the central straight line, pointing + to -
  \draw[engblue, ->] (-0.2,0.001) -- (0.2,0.001);
  \filldraw[engred] (P) circle [radius=4pt];
  \node[white, font=\scriptsize\bfseries] at (P) {$+$};
  \filldraw[engblue] (N) circle [radius=4pt];
  \node[white, font=\scriptsize\bfseries] at (N) {$-$};
  \node[font=\small] at (0,-2.2) {dipole ($+q$ and $-q$)};
\end{scope}
\end{tikzpicture}
\caption[Field lines of a point charge and a dipole]{Electric field
lines for an isolated positive point charge (left) and for an
equal-and-opposite charge pair, the dipole (right). Lines begin on
positive charge and end on negative charge, the geometric statement of
Gauss's law: the net outward flux through any surface counts the charge
enclosed. For the isolated charge the lines run radially to infinity
and the field falls as $1/r^{2}$; for the dipole the lines close from
the positive charge onto the negative charge and the far field falls
faster, as $1/r^{3}$.}
\label{fig:vol04ch08:point-charge-field}
\end{figure}


The combined force is the \engterm{Lorentz force},
\[
\vect{F} = q(\vect{E} + \vect{v}\times \vect{B}),
\]
the empirical input that closes the loop between fields and motion.
Once we know the fields, the Lorentz force tells us what charges do;
once we know what the charges do, Maxwell's equations tell us what the
fields do. The two halves close.

\begin{archetype}[Balance, transport, interface]
Every electromagnetic equation in the chapter is one of three patterns.
A \engterm{balance} equates a flux across a surface to the source
enclosed (Gauss's law, the no-monopole law). A \engterm{transport}
equation relates the time rate of change of one field to the curl of
another (Faraday, Amp\`ere-Maxwell). An \engterm{interface} condition
relates the field on one side of a material boundary to the field on
the other (boundary conditions on $\vect{E}$, $\vect{D}$, $\vect{B}$,
$\vect{H}$). The discipline of the chapter is to recognise which
pattern a problem requires before reaching for the formula.
\end{archetype}

\section{Gauss's law for electricity}\pathtag{core}

The first Maxwell equation, in integral form, is
\[
\oint_{\partial V} \vect{E}\cdot \dd\vect{A} = \frac{Q_{\text{enc}}}{
\varepsilon_{0}},
\]
where $\partial V$ is a closed surface bounding the volume $V$,
$Q_{\text{enc}} = \int_{V}\rho\,\dd V$ is the charge enclosed, and
$\varepsilon_{0} = 8.854\,187\,\times 10^{-12}\,\si{\farad\per\meter}$
is the permittivity of free space. The law is the experimental content
of Coulomb's $1/r^{2}$ result, extended from point charges to arbitrary
distributions: the total outward flux of $\vect{E}$ through a closed
surface depends only on the charge enclosed.

The differential form follows from the divergence theorem of
Volume II, Chapter 8:
\[
\divg \vect{E} = \frac{\rho}{\varepsilon_{0}}.
\]
The electric field's divergence at a point equals the charge density at
that point, scaled by $1/\varepsilon_{0}$. In a region with no charge
the divergence is zero; field lines neither begin nor end. In a region
with positive charge density, field lines diverge from the charge; with
negative charge density, field lines converge.

\begin{principle}[Choose the symmetry first]
Gauss's law gives the field by quadrature only when the symmetry of the
problem makes $\vect{E}\cdot\dd\vect{A}$ constant on a chosen Gaussian
surface. Three symmetries close the integral:
\begin{itemize}[leftmargin=*]
  \item Spherical (point charge, charged sphere): Gaussian sphere
    centred on the source.
  \item Cylindrical (long wire, charged cylinder): Gaussian cylinder
    coaxial with the source.
  \item Planar (infinite sheet, parallel-plate capacitor):
    Gaussian pillbox spanning the sheet.
\end{itemize}
If none of the three applies, Gauss's law is still true but no longer
a tool for computing the field; the field is obtained instead by
integrating Coulomb's law or by solving Poisson's equation. The
engineer's instinct is to identify the symmetry before reaching for
the integral.
\end{principle}

% Figure: the three symmetric Gaussian surfaces that close Gauss's law.
% Spherical surface around a point charge; cylindrical surface around a
% line charge; pillbox straddling a charged sheet.
\begin{figure}[htbp]
\centering
\begin{tikzpicture}[>=Latex, scale=1.0]
% === Sphere around point charge ===
\begin{scope}[shift={(0,0)}]
  \draw[engblue, thick] (0,0) circle [radius=1.3];
  \draw[engblue, thick, dashed] (0,0) ellipse [x radius=1.3, y radius=0.45];
  \filldraw[engred] (0,0) circle [radius=3pt];
  \node[white, font=\tiny\bfseries] at (0,0) {$+$};
  \draw[enggray, ->] (0,0) -- (35:1.3) node[midway, above, font=\tiny] {$r$};
  \node[font=\small] at (0,-2.0) {spherical: point charge};
\end{scope}
% === Cylinder around line charge ===
\begin{scope}[shift={(5.0,0)}]
  \draw[engred, very thick] (0,-1.6) -- (0,1.6);
  \node[engred, font=\tiny] at (0.25,1.5) {$\lambda$};
  \draw[engblue, thick] (-0.8,-1.1) -- (-0.8,1.1);
  \draw[engblue, thick] (0.8,-1.1) -- (0.8,1.1);
  \draw[engblue, thick] (0,1.1) ellipse [x radius=0.8, y radius=0.28];
  \draw[engblue, thick, dashed] (0,-1.1) ellipse [x radius=0.8, y radius=0.28];
  \draw[enggray, ->] (0,-0.3) -- (0.8,-0.3) node[midway, below, font=\tiny] {$r$};
  \node[font=\small] at (0,-2.0) {cylindrical: line charge};
\end{scope}
% === Pillbox across a charged sheet ===
\begin{scope}[shift={(10.0,0)}]
  \fill[engred!25] (-1.4,-0.05) rectangle (1.4,0.05);
  \draw[engred] (-1.4,0) -- (1.4,0);
  \node[engred, font=\tiny] at (1.55,0.18) {$\sigma$};
  \draw[engblue, thick] (-0.6,-0.7) rectangle (0.6,0.7);
  \draw[engblue, ->] (0,0.7) -- (0,1.2) node[above, font=\tiny] {$\vect{E}$};
  \draw[engblue, ->] (0,-0.7) -- (0,-1.2) node[below, font=\tiny] {$\vect{E}$};
  \node[font=\small] at (0,-2.0) {planar: charged sheet};
\end{scope}
\end{tikzpicture}
\caption[The three symmetric Gaussian surfaces]{The three symmetries
that let Gauss's law deliver the field by quadrature. A spherical
surface centred on a point or spherical charge (left) gives
$E = q/(4\pi\varepsilon_0 r^2)$; a coaxial cylinder around a long line
charge (centre) gives $E = \lambda/(2\pi\varepsilon_0 r)$; a pillbox
straddling an infinite charged sheet (right) gives
$E = \sigma/(2\varepsilon_0)$ on each face. When none of the three fits,
Gauss's law remains true but is no longer a tool for computing the
field.}
\label{fig:vol04ch08:gaussian-surfaces}
\end{figure}


A point charge $q$ at the origin produces, by spherical symmetry on
a Gaussian sphere of radius $r$,
\[
4\pi r^{2}\, E(r) = \frac{q}{\varepsilon_{0}}, \quad E(r) =
\frac{q}{4\pi \varepsilon_{0} r^{2}},
\]
which is Coulomb's law recovered. The constant $1/(4\pi\varepsilon_{0})
\approx 8.988\times 10^{9}\,\si{\newton\meter\squared\per\coulomb
\squared}$ is large; even small charges produce large forces at small
separations, which is why static-electricity sparks are uncomfortable
and why most matter is electrically neutral at the macroscopic scale.

An infinite line charge with linear density $\lambda$ in $\si{\coulomb
\per\meter}$ produces, by cylindrical symmetry on a Gaussian cylinder
of radius $r$ and length $L$,
\[
2\pi r L\, E(r) = \frac{\lambda L}{\varepsilon_{0}}, \quad E(r) =
\frac{\lambda}{2\pi \varepsilon_{0} r},
\]
falling as $1/r$ rather than $1/r^{2}$. Inside a coaxial cable the
field between conductors has this form, and the result drives the
cable's capacitance per unit length, derived in Section 8.8.

An infinite sheet with surface charge density $\sigma$ in $\si{\coulomb
\per\meter\squared}$ produces, by planar symmetry on a pillbox
straddling the sheet,
\[
E = \frac{\sigma}{2\varepsilon_{0}},
\]
on each side, normal to the sheet, independent of distance. A pair of
infinite parallel sheets with equal and opposite charge densities
$\pm\sigma$ produces $E = \sigma/\varepsilon_{0}$ between them and zero
outside; this is the working model of the ideal parallel-plate
capacitor.

\section{No magnetic monopoles}\pathtag{core}

The second Maxwell equation states that the magnetic flux through any
closed surface is zero,
\[
\oint_{\partial V} \vect{B}\cdot \dd\vect{A} = 0,
\]
or equivalently
\[
\divg \vect{B} = 0.
\]
The empirical content is that no magnetic monopole has ever been
observed despite repeated searches in the 20th and 21st centuries; the
upper bound on cosmic monopole flux is so small that for engineering
purposes monopoles do not exist. Every magnetic field line that enters
a closed surface leaves it; magnetic field lines form closed loops.

A consequence is that the magnetic field can be written as the curl of
a vector potential,
\[
\vect{B} = \curl \vect{A},
\]
since the divergence of any curl is zero. The vector potential
$\vect{A}$ is not unique: adding the gradient of any scalar function
leaves $\vect{B}$ unchanged, a freedom called \engterm{gauge
invariance}. The freedom is fixed by a gauge condition; the Coulomb
gauge $\divg \vect{A} = 0$ and the Lorenz gauge $\divg \vect{A} +
\mu_{0}\varepsilon_{0}\,\partial \phi/\partial t = 0$ are the two
most useful.

The no-monopole equation does not directly give a field by quadrature
the way Gauss's law does, but it constrains the topology of every
magnetic field the engineer will encounter. A solenoid's field exits
one end and re-enters the other in a closed-loop pattern; a bar
magnet's field lines emerge from the north pole and return to the
south through space outside and through the magnet inside; a current
loop's field circulates around the loop's axis. The closed-loop
geometry is not optional; it is the consequence of the empirical
absence of monopoles.

% Figure: solenoid with interior axial field and exterior return path
% forming closed loops (no-monopole law). Schematic windings shown as
% circles, interior field uniform, exterior weak.
\begin{figure}[htbp]
\centering
\begin{tikzpicture}[>=Latex, scale=1.0]
% winding turns (cross sections): top row dots (current out), bottom crosses (in)
\foreach \x in {-2.4,-1.8,...,2.4}{
  \draw[engred, thick] (\x,0.9) circle [radius=0.12];
  \filldraw[engred] (\x,0.9) circle [radius=1.6pt];        % current out of page
  \draw[engred, thick] (\x,-0.9) circle [radius=0.12];
  \draw[engred, thick] ({\x-0.085},{-0.985}) -- ({\x+0.085},{-0.815});  % cross
  \draw[engred, thick] ({\x-0.085},{-0.815}) -- ({\x+0.085},{-0.985});
}
% interior uniform field
\foreach \y in {-0.5,-0.25,0,0.25,0.5}{
  \draw[engblue, ->] (-2.2,\y) -- (2.2,\y);
}
\node[engblue, font=\small] at (0,0.72) {$B=\mu_0 n I$ (uniform inside)};
% exterior closed return loops
\draw[enggray, ->] (2.4,0.55) .. controls (3.6,0.4) and (3.6,-0.4) .. (2.4,-0.55);
\draw[enggray] (-2.4,0.55) .. controls (-3.6,0.4) and (-3.6,-0.4) .. (-2.4,-0.55);
\draw[enggray, ->] (-2.4,-0.55) -- (2.4,-0.55) coordinate[pos=0.5];
\draw[enggray, ->] (2.4,0.55) -- (-2.4,0.55);
\node[enggray, font=\scriptsize] at (3.5,0) {weak return field};
\node[font=\small] at (0,-1.6) {$n$ turns per metre, current $I$};
\end{tikzpicture}
\caption[Solenoid field and closed loops]{A long solenoid carries $n$
turns per unit length at current $I$. The interior field is uniform and
axial, $B = \mu_0 n I$, while the exterior field is weak and spread over
the return path. Every interior line that leaves one end re-enters the
other: magnetic field lines close on themselves, the geometric content
of $\divg\vect{B}=0$ and the empirical absence of magnetic monopoles.
The dots mark current out of the page on the top winding, the crosses
current into the page on the bottom.}
\label{fig:vol04ch08:solenoid-field}
\end{figure}


\section{Faraday's law}\pathtag{core}

The third Maxwell equation is Faraday's law of induction:
\[
\oint_{C} \vect{E}\cdot \dd\vect{\ell} = -\dv{}{t}\int_{S} \vect{B}
\cdot \dd\vect{A},
\]
where $C$ is a closed curve bounding the surface $S$, and the time
derivative acts on the total magnetic flux through $S$. The line
integral on the left is the \engterm{electromotive force} (EMF)
around $C$; the surface integral on the right is the magnetic flux.
The differential form,
\[
\curl \vect{E} = -\pdv{\vect{B}}{t},
\]
follows from Stokes's theorem. A changing magnetic field produces a
circulating electric field; the minus sign (Lenz's law) is the
empirical content that the induced EMF opposes the change.

% Figure: Faraday induction. A bar magnet approaching a loop, the changing
% flux, and the induced current whose own field opposes the change (Lenz).
\begin{figure}[htbp]
\centering
\begin{tikzpicture}[>=Latex, scale=1.0]
% the loop drawn as an ellipse seen at an angle
\draw[engblue, very thick] (0,0) ellipse [x radius=1.0, y radius=1.5];
% induced current arrow on the loop
\draw[engblue, ->] (1.0,0.4) arc [start angle=20, end angle=70, x radius=1.0, y radius=1.5];
\node[engblue, font=\scriptsize] at (1.5,1.1) {$I_{\text{ind}}$};
% bar magnet approaching from the left
\fill[engred!70] (-4.0,-0.35) rectangle (-2.6,0.35);
\fill[engblue!50] (-2.6,-0.35) rectangle (-2.2,0.35);
\node[white, font=\scriptsize\bfseries] at (-3.3,0) {N};
\node[white, font=\scriptsize\bfseries] at (-2.4,0) {S};
\draw[->, thick] (-2.0,0.9) -- (-1.0,0.9) node[midway, above, font=\scriptsize] {motion};
% magnet field lines threading the loop
\draw[enggray, ->] (-2.2,0.15) -- (-0.3,0.15);
\draw[enggray, ->] (-2.2,-0.15) -- (-0.3,-0.15);
\node[enggray, font=\scriptsize] at (-1.2,0.42) {$\Phi_B$ increasing};
% EMF annotation
\node[font=\small] at (0,-2.2) {$\mathcal{E}=-N\,\dd\Phi_B/\dd t$};
\end{tikzpicture}
\caption[Faraday induction and Lenz's law]{A magnet moved toward a
conducting loop increases the magnetic flux $\Phi_B$ through the loop.
Faraday's law drives an electromotive force $\mathcal{E} = -N\,
\dd\Phi_B/\dd t$ around the loop, and the induced current $I_{\text{ind}}$
flows in the direction whose own magnetic field opposes the increase
(Lenz's law, the minus sign). The opposing field produces a force that
resists the magnet's approach: the back-EMF that brakes a generator and
recovers kinetic energy in regenerative braking. Without the minus sign
energy would not balance.}
\label{fig:vol04ch08:faraday-loop}
\end{figure}


Faraday's law is the working principle of every generator, transformer,
and inductive coupling in engineering. The flux through a coil of $N$
turns is $N\Phi_{B}$ where $\Phi_{B} = \int \vect{B}\cdot \dd\vect{A}$
through a single turn; the induced EMF is
\[
\mathcal{E} = -N\dv{\Phi_{B}}{t}.
\]
A transformer with primary voltage $V_{1}$ and turns $N_{1}$ produces
on a secondary with $N_{2}$ turns
\[
V_{2} = V_{1}\frac{N_{2}}{N_{1}},
\]
in the ideal limit of zero leakage flux and zero winding resistance.
The transformer ratio is the source-text reason the grid uses
high-voltage transmission: stepping voltage up at the generator and
down at the consumer reduces $I^{2}R$ losses by the factor $(N_{1}/
N_{2})^{2}$.

The minus sign carries operational consequences. A coil moved into a
field experiences an EMF that drives current in the direction that
\emph{opposes} the motion; the coil feels a retarding force. This is
the principle of regenerative braking in electric vehicles: kinetic
energy is converted to electrical energy by the back-EMF of the
motor's coils, with the mechanical force decelerating the vehicle.
Without the minus sign, kinetic energy could grow without bound from
nothing; the minus sign is conservation of energy in disguise.

\begin{estimation}
\textbf{Question.} A bicycle dynamo lights a $3\,\si{\watt}$ headlamp
when the wheel turns at $20\,\si{\kilo\meter\per\hour}$. Estimate the
magnetic flux change per revolution of the dynamo's rotor and the
peak voltage induced.

\textbf{Estimate, before reading on.} The reader is asked to state a
factor-of-three-confidence answer before reading the working.

\textbf{Working.} The dynamo rotates against the tyre; at $20\,\si{
\kilo\meter\per\hour}$ the wheel's circumference of about $2\,\si{
\meter}$ rotates roughly $3\,\si{\hertz}$. A typical dynamo has a
gear ratio of around $5$ to $10$ between tyre and rotor, giving a
rotor frequency of $15$ to $30\,\si{\hertz}$. The lamp is $3\,\si{
\watt}$ at $6\,\si{\volt}$, so RMS current is $0.5\,\si{\ampere}$ and
peak voltage in a single sinusoid would be about $8.5\,\si{\volt}$.
The peak voltage occurs at maximum $\dd\Phi/\dd t$; with $N$ turns
and peak flux $\Phi_{0}$ in a rotor turning at $\omega$, peak EMF is
$N\Phi_{0}\omega$. For $\omega \approx 2\pi\times 20 = 125\,\si{\per
\second}$ and $V_{\text{peak}} \approx 8.5\,\si{\volt}$, the product
$N\Phi_{0} \approx 0.07\,\si{\weber}$. A coil of $200$ turns then needs
$\Phi_{0} \approx 350\,\si{\micro\weber}$, consistent with a $1\,\si{
\centi\meter\squared}$ pole face seeing $\approx 0.35\,\si{\tesla}$
from a small ferrite magnet.

\textbf{Postmortem.} The estimate hangs on the gear ratio and the
number of turns; both can be observed on a real dynamo. A reader who
guessed within a factor of three internalised the order of magnitude
of generator design without writing a single equation in coordinate
form. The empirical numbers are dyamos sold for bicycles, current as
of 2025.
\end{estimation}

\section{Amp\`ere-Maxwell law}\pathtag{core}

The fourth Maxwell equation, in the form Maxwell himself completed it,
is
\[
\oint_{C} \vect{B}\cdot \dd\vect{\ell} = \mu_{0}\left( I_{\text{enc}}
+ \varepsilon_{0}\dv{}{t}\int_{S}\vect{E}\cdot \dd\vect{A}\right),
\]
or equivalently in differential form
\[
\curl \vect{B} = \mu_{0}\vect{J} + \mu_{0}\varepsilon_{0}\pdv{\vect{E}
}{t},
\]
where $\mu_{0} = 4\pi\times 10^{-7}\,\si{\henry\per\meter}$ is the
permeability of free space (its value redefined by the 2019 SI as
an experimentally determined quantity, no longer exact). The first
term on the right is the current density; the second term is Maxwell's
\engterm{displacement current}, the time rate of change of the
electric field.

The displacement current is the term Maxwell added in 1865 to make the
four equations consistent with charge conservation. Without it,
$\divg \curl \vect{B} = \mu_{0}\divg \vect{J} = -\mu_{0}\partial \rho/
\partial t$, which is non-zero whenever charge accumulates anywhere;
but $\divg \curl$ is identically zero for any field, so the equation
is inconsistent without the displacement current. With the
displacement current,
\[
\divg(\mu_{0}\vect{J} + \mu_{0}\varepsilon_{0}\pdv{\vect{E}}{t}) =
\mu_{0}(-\pdv{\rho}{t}) + \mu_{0}\varepsilon_{0}\pdv{(\rho/
\varepsilon_{0})}{t} = 0,
\]
identically, by continuity. The mathematical requirement of consistency
forced the physical prediction that a changing electric field produces
a magnetic field even in the absence of current, and the prediction
was confirmed empirically by Hertz in 1888 \cite{hist:hertz1893}.

% Figure: displacement current in a charging capacitor. Amperian loop
% around the wire; two candidate surfaces (one cut by the wire, one
% passing between the plates) must give the same enclosed "current".
\begin{figure}[htbp]
\centering
\begin{tikzpicture}[>=Latex, scale=1.0]
% wire from left into the left plate
\draw[engred, line width=1.5pt] (-4.2,0) -- (-1.2,0);
\node[engred, font=\scriptsize, anchor=south] at (-3.5,0.05) {$I$};
\draw[engred, ->] (-3.0,0) -- (-2.6,0);
% capacitor plates
\draw[engblue, line width=2pt] (-1.2,-0.9) -- (-1.2,0.9);
\draw[engblue, line width=2pt] (1.2,-0.9) -- (1.2,0.9);
% wire out the right plate
\draw[engred, line width=1.5pt] (1.2,0) -- (4.2,0);
\draw[engred, ->] (2.6,0) -- (3.0,0);
% changing E field between plates
\foreach \y in {-0.55,-0.275,0,0.275,0.55}{
  \draw[enggray, ->] (-1.1,\y) -- (1.1,\y);
}
\node[font=\scriptsize] at (0,0.75) {$\partial\vect{E}/\partial t$};
% Amperian loop (around the right wire), shown as a circle in the plane
\draw[engamber, thick] (-2.6,0) ellipse [x radius=0.35, y radius=0.95];
\node[engamber, font=\scriptsize, anchor=east] at (-2.95,0.9) {Amperian};
\node[engamber, font=\scriptsize, anchor=east] at (-2.95,0.6) {loop $C$};
% surface S1: a flat disk cut by the wire
\node[font=\scriptsize, anchor=north] at (-2.6,-1.0) {$S_1$: cut by $I$};
% surface S2: bulges between the plates, no conduction current, displacement current
\draw[engamber, dashed] (-2.6,0.95) .. controls (0,2.0) and (0,-2.0) .. (-2.6,-0.95);
\node[font=\scriptsize, anchor=south] at (0,1.5) {$S_2$: cut by $\varepsilon_0\,\partial\vect{E}/\partial t$};
\end{tikzpicture}
\caption[Displacement current in a charging capacitor]{Why Maxwell's
displacement current is needed. The Amperian loop $C$ bounds infinitely
many surfaces. The flat surface $S_1$ is pierced by the conduction
current $I$ in the wire; the bulging surface $S_2$ passes between the
capacitor plates, where no charge crosses, yet both must give the same
result for $\oint\vect{B}\cdot\dd\vect{\ell}$. The changing electric
field between the plates supplies the missing term: the displacement
current $\varepsilon_0\,\dd/\dd t\int\vect{E}\cdot\dd\vect{A}$ through
$S_2$ equals the conduction current $I$ through $S_1$. Consistency with
charge conservation forced the term, and Hertz confirmed its
consequence, electromagnetic waves, in 1888.}
\label{fig:vol04ch08:displacement-current}
\end{figure}


For a long straight wire carrying steady current $I$, Amp\`ere's law on
a circle of radius $r$ centred on the wire gives, by cylindrical
symmetry,
\[
2\pi r B(r) = \mu_{0} I, \quad B(r) = \frac{\mu_{0} I}{2\pi r}.
\]
A wire carrying $1\,\si{\ampere}$ at $1\,\si{\centi\meter}$ from its
axis produces a field of $20\,\si{\micro\tesla}$, comparable to the
Earth's field; the compass-deflection experiment that Oersted
performed in 1820 is recoverable on a bench with no special equipment.

A long solenoid with $n$ turns per unit length carrying current $I$
produces an axial field
\[
B = \mu_{0} n I
\]
inside, uniform across the cross-section in the limit of an infinite
solenoid; outside, the field is negligible compared to the interior.
A solenoid with $n = 1000\,\si{\per\meter}$ at $I = 10\,\si{\ampere}$
produces $B = 12.6\,\si{\milli\tesla}$, a useful field for laboratory
purposes.

% Figure: magnetic field of a long straight current-carrying wire. The 1/r
% decay plotted, with the circular-loop geometry inset as an annotation.
\begin{figure}[htbp]
\centering
\begin{tikzpicture}
\begin{axis}[
  width=12cm, height=7cm,
  xlabel={radial distance from wire, $r$ (cm)},
  ylabel={field magnitude $B$ ($\mu$T)},
  xmin=0.5, xmax=20, ymin=0, ymax=42,
  axis lines=left,
  tick label style={font=\scriptsize},
  label style={font=\small},
  legend pos=north east,
  legend style={font=\scriptsize},
]
% B = mu0 I / (2 pi r), I = 1 A: B(uT) = 0.2 / r(cm) * 100 = 20/r(cm)
\addplot[engblue, very thick, samples=120, domain=0.5:20] {20/x};
\addlegendentry{$I=1$\,A}
% I = 0.5 A
\addplot[engred, thick, dashed, samples=120, domain=0.5:20] {10/x};
\addlegendentry{$I=0.5$\,A}
% mark the 1 cm, 20 uT point
\filldraw[engblue] (axis cs:1,20) circle [radius=1.6pt];
\node[engblue, font=\scriptsize, anchor=west] at (axis cs:1.4,21) {$20\,\mu$T at $1$\,cm};
\end{axis}
\end{tikzpicture}
\caption[Magnetic field of a straight wire]{Magnetic field magnitude
around a long straight wire, $B = \mu_0 I/(2\pi r)$, the result Amp\`ere's
law gives by cylindrical symmetry. At $1\,\si{\centi\meter}$ from a wire
carrying $1\,\si{\ampere}$ the field is about $20\,\si{\micro\tesla}$,
comparable to the Earth's field, which is why Oersted's 1820 compass
deflection is recoverable on any bench. The $1/r$ decay is the law the
chapter project sets out to verify by measurement.}
\label{fig:vol04ch08:wire-bfield}
\end{figure}


The same Biot-Savart integral that yields these symmetric closed forms
also handles geometries without symmetry, where the only route is direct
quadrature. The script \repopath{code/biot\_savart\_segment.py} integrates
$\dd\vect{B} = (\mu_0 I/4\pi)\,\dd\vect{\ell}\times\unitvec{r}/r^{2}$
along a finite straight segment and along the four sides of a square
loop; it recovers the infinite-wire result $\mu_0 I/(2\pi r)$ as the
segment lengthens and the centre-of-loop value $2\sqrt{2}\,\mu_0 I/
(\pi a)$ for a square of side $a$. The lesson of the symmetric cases is
that they are the exceptions; most real coils are integrated
numerically.

% Figure: the four Maxwell equations as a summary table-diagram, paired with
% their one-line physical reading and the archetype each belongs to.
\begin{figure}[htbp]
\centering
\begin{tikzpicture}[font=\small]
\tikzset{
  eqbox/.style={draw=engblue, thick, rounded corners, fill=engblue!6,
    text width=4.6cm, align=center, inner sep=6pt, minimum height=1.5cm},
  readbox/.style={draw=enggray, rounded corners, fill=enggray!8,
    text width=5.2cm, align=left, inner sep=6pt, minimum height=1.5cm},
}
% rows
\node[eqbox] (g1) at (0,4.8) {Gauss (electric)\\[2pt] $\displaystyle\divg\vect{E}=\rho/\varepsilon_0$};
\node[readbox, right=0.6cm of g1] (r1) {Charge is the source of $\vect{E}$. Flux counts enclosed charge. \emph{Balance.}};
\node[eqbox] (g2) at (0,3.0) {Gauss (magnetic)\\[2pt] $\displaystyle\divg\vect{B}=0$};
\node[readbox, right=0.6cm of g2] (r2) {No magnetic charge. $\vect{B}$ lines close on themselves. \emph{Balance.}};
\node[eqbox] (g3) at (0,1.2) {Faraday\\[2pt] $\displaystyle\curl\vect{E}=-\partial\vect{B}/\partial t$};
\node[readbox, right=0.6cm of g3] (r3) {Changing $\vect{B}$ drives a circulating $\vect{E}$. \emph{Transport.}};
\node[eqbox] (g4) at (0,-0.6) {Amp\`ere--Maxwell\\[2pt] $\displaystyle\curl\vect{B}=\mu_0\vect{J}+\mu_0\varepsilon_0\,\partial\vect{E}/\partial t$};
\node[readbox, right=0.6cm of g4] (r4) {Current and changing $\vect{E}$ drive a circulating $\vect{B}$. \emph{Transport.}};
\end{tikzpicture}
\caption[The four Maxwell equations]{The four Maxwell equations in
differential form, each with its physical reading and its archetype.
The two divergence laws are balances: they equate a flux to its enclosed
source, and the magnetic one has no source at all. The two curl laws are
transport equations: each ties the time rate of change of one field to
the circulation of the other, and together they imply the wave equation
and the finite speed $c = 1/\sqrt{\mu_0\varepsilon_0}$. The discipline
of the chapter is to recognise which archetype a problem needs before
reaching for a formula.}
\label{fig:vol04ch08:maxwell-summary}
\end{figure}


With the four equations stated, their differential forms collected in
\cref{fig:vol04ch08:maxwell-summary} read as two balances and two
transport laws. The remaining sections draw the consequences: the energy
the fields carry, the conditions they obey at material boundaries, and
the failure modes they produce in real systems.

\section{Electromagnetic energy and the Poynting vector}\pathtag{core}

The energy density stored in electric and magnetic fields is
\[
u_{e} = \tfrac{1}{2}\varepsilon_{0}|\vect{E}|^{2}, \quad u_{m} =
\tfrac{1}{2\mu_{0}}|\vect{B}|^{2},
\]
both in $\si{\joule\per\meter\cubed}$. A parallel-plate capacitor at
$1\,\si{\kilo\volt}$ with $1\,\si{\milli\meter}$ plate spacing has
$E = 10^{6}\,\si{\volt\per\meter}$ and $u_{e} = 4.4\,\si{\joule\per
\meter\cubed}$; the energy in a $1\,\si{\centi\meter\cubed}$ volume is
about $4\,\si{\micro\joule}$. A MRI magnet at $3\,\si{\tesla}$ has
$u_{m} = 3.6\,\si{\mega\joule\per\meter\cubed}$; the energy stored in a
$1\,\si{\meter\cubed}$ bore is $3.6\,\si{\mega\joule}$, comparable to
the kinetic energy of a small car at $100\,\si{\kilo\meter\per\hour}$.
The magnetic field is the heavier energy reservoir for the same volume
once $B$ exceeds the order of $\sqrt{\mu_{0}\varepsilon_{0}}\,E$,
which is to say $E/B \sim c$, the speed of light.

The rate at which electromagnetic energy flows through a surface is the
flux of the \engterm{Poynting vector},
\[
\vect{S} = \frac{1}{\mu_{0}}\vect{E}\times \vect{B},
\]
in $\si{\watt\per\meter\squared}$. The Poynting theorem, derived by
taking $\vect{E}\cdot$ (the Amp\`ere-Maxwell equation) minus
$\vect{B}\cdot$ (Faraday's equation) and integrating, states
\[
\pdv{(u_{e}+u_{m})}{t} + \divg \vect{S} = -\vect{J}\cdot \vect{E},
\]
the local energy balance: the rate of change of stored field energy
plus the divergence of the Poynting vector equals the negative of the
work done by the field on the current. The right-hand side is the
mechanism by which field energy becomes thermal energy in a resistor:
$\vect{J}\cdot \vect{E}$ is the volumetric power density dissipated.

\begin{principle}[Energy travels in the fields, not in the wire]
For a coaxial cable carrying DC current from a source to a load, the
electric field points radially inward, the magnetic field circulates
azimuthally, and the Poynting vector $\vect{E}\times \vect{B}/\mu_{0}$
points along the cable's axis from source to load. Integrating over
the cable's cross-section gives a power flux equal to $VI$, the
familiar product of voltage and current. The energy is carried in the
space between the conductors, not in the conductors themselves. This
is more than a curiosity: it is the working content of transmission-
line theory in Chapter 10, where impedance and standing waves emerge
from the field's path between conductors.
\end{principle}

A plane electromagnetic wave in vacuum has $|\vect{E}|/|\vect{B}| = c
= 1/\sqrt{\mu_{0}\varepsilon_{0}} = 2.998\times 10^{8}\,\si{\meter\per
\second}$, and equal contributions to $u_{e}$ and $u_{m}$. The
time-averaged Poynting vector magnitude for a sinusoidal plane wave
of peak electric-field amplitude $E_{0}$ is
\[
\langle S\rangle = \frac{1}{2}\frac{E_{0}^{2}}{\mu_{0} c} =
\frac{1}{2}\varepsilon_{0} c E_{0}^{2}.
\]
Solar irradiance at the top of Earth's atmosphere is about $1361\,\si{
\watt\per\meter\squared}$ (the \engterm{solar constant}, current as
of 2025); equating this to $\langle S\rangle$ gives a peak field of
$E_{0} \approx 1000\,\si{\volt\per\meter}$. The same calculation drives
the design of radio receivers (Chapter 10) and photovoltaic cells
(Chapter 12). The script \repopath{code/plane\_wave.py} carries the
conversion both ways, from irradiance to peak field amplitudes and back,
together with the time-averaged energy density and the radiation
pressure on an absorbing surface; at the solar constant it returns
$E_{0} = 1013\,\si{\volt\per\meter}$, $B_{0} = 3.4\,\si{\micro\tesla}$,
and a radiation pressure of $4.5\,\si{\micro\pascal}$.

% Figure: a plane electromagnetic wave. E in the x-z plane, B in the y-z
% plane, propagation along z, E perpendicular to B perpendicular to k.
\begin{figure}[htbp]
\centering
\begin{tikzpicture}
\begin{axis}[
  width=13cm, height=6.5cm,
  xlabel={propagation direction $z$},
  axis lines=middle,
  xmin=0, xmax=12.6, ymin=-1.4, ymax=1.4,
  xtick=\empty, ytick=\empty,
  clip=false,
  label style={font=\small},
]
% E field: sine in the vertical plane
\addplot[engred, very thick, samples=160, domain=0:12.566] {sin(deg(x))};
\node[engred, font=\small] at (axis cs:1.7,1.25) {$\vect{E}$};
% B field: cosine? in phase with E for a wave, drawn as a phase-locked sine
% scaled to look like the transverse companion (shown smaller for the page)
\addplot[engblue, very thick, samples=160, domain=0:12.566] {0.55*sin(deg(x))};
\node[engblue, font=\small] at (axis cs:1.7,0.78) {$\vect{B}$};
% propagation arrow
\draw[->, thick] (axis cs:11.0,0) -- (axis cs:12.3,0) node[above, font=\small] {$c$};
\end{axis}
\end{tikzpicture}
\caption[Plane electromagnetic wave]{A plane electromagnetic wave in
vacuum. The electric field $\vect{E}$ and magnetic field $\vect{B}$
oscillate in phase, transverse to each other and to the propagation
direction, with $|\vect{E}|/|\vect{B}| = c = 1/\sqrt{\mu_0\varepsilon_0}$.
The two curl equations of Maxwell, Faraday and Amp\`ere--Maxwell, feed
each other: a changing $\vect{E}$ sources $\vect{B}$ and a changing
$\vect{B}$ sources $\vect{E}$, and the pair propagates at $c$. The energy
density splits equally between the two fields, and the time-averaged
Poynting vector $\langle S\rangle = \tfrac12\varepsilon_0 c E_0^2$ is the
intensity the wave delivers. The two curves are drawn at different
heights only for legibility; their amplitudes are fixed by the ratio
$c$.}
\label{fig:vol04ch08:em-wave}
\end{figure}


\begin{mastery}
The wave itself falls out of the two curl equations in vacuum, where
$\rho = 0$ and $\vect{J} = \vect{0}$. Take the curl of Faraday's law,
$\curl(\curl\vect{E}) = -\partial(\curl\vect{B})/\partial t$, and
substitute the Amp\`ere-Maxwell law $\curl\vect{B} = \mu_0\varepsilon_0\,
\partial\vect{E}/\partial t$ on the right. The vector identity
$\curl(\curl\vect{E}) = \grad(\divg\vect{E}) - \laplacian\vect{E}$ with
$\divg\vect{E} = 0$ collapses the left side, leaving
\[
\laplacian\vect{E} = \mu_0\varepsilon_0\frac{\partial^{2}\vect{E}}{
\partial t^{2}},
\]
the wave equation with speed $c = 1/\sqrt{\mu_0\varepsilon_0}$. The
identical manipulation on $\vect{B}$ gives the same equation. The speed
came out of two laboratory constants measured in electrostatics and
magnetostatics, and it matched the measured speed of light; that
coincidence is the content of the chapter's epigraph. The same
second-order wave operator governed the vibrating string in Volume II
and recurs in transmission-line theory in Chapter 10.
\end{mastery}

\section{Boundary conditions and material laws}\pathtag{standard}

Maxwell's equations in their bare form describe fields in vacuum. In
material media, the bound charges and bound currents of the material
modify the fields. The book's working notation introduces two
auxiliary fields: the \engterm{electric displacement} $\vect{D}$ and
the \engterm{magnetic field intensity} $\vect{H}$, defined by
\[
\vect{D} = \varepsilon_{0}\vect{E} + \vect{P}, \quad \vect{H} =
\frac{1}{\mu_{0}}\vect{B} - \vect{M},
\]
where $\vect{P}$ is the polarisation density (dipole moment per unit
volume) and $\vect{M}$ is the magnetisation (magnetic moment per unit
volume). The Maxwell equations in matter become
\[
\divg \vect{D} = \rho_{f}, \quad \divg \vect{B} = 0,
\]
\[
\curl \vect{E} = -\pdv{\vect{B}}{t}, \quad \curl \vect{H} = \vect{J}_{f}
+ \pdv{\vect{D}}{t},
\]
where the subscript $f$ denotes free (as opposed to bound) charges and
currents. For linear isotropic materials,
\[
\vect{D} = \varepsilon \vect{E}, \quad \vect{B} = \mu \vect{H},
\quad \vect{J}_{f} = \sigma \vect{E},
\]
introducing the \engterm{permittivity} $\varepsilon$, the
\engterm{permeability} $\mu$, and the \engterm{conductivity} $\sigma$.
The relative quantities are $\varepsilon_{r} = \varepsilon/\varepsilon
_{0}$ and $\mu_{r} = \mu/\mu_{0}$. Water has $\varepsilon_{r} \approx
80$ at low frequencies; ferrite has $\mu_{r}$ from a few to $10\,000$;
copper has $\sigma \approx 5.96\times 10^{7}\,\si{\siemens\per\meter}$;
glass at room temperature has $\sigma \approx 10^{-12}\,\si{\siemens
\per\meter}$.

The third Maxwell-in-matter equation is Ohm's law in field form:
$\vect{J} = \sigma \vect{E}$ is the constitutive equation of a linear
conductor, and the lumped-element relation $V = IR$ derived in
Chapter 9 is its integral over a wire's cross-section and length.

At a sharp boundary between materials $1$ and $2$ with unit normal
$\unitvec{n}$ from $1$ to $2$, the field components obey:
\begin{itemize}[leftmargin=*]
  \item $D_{2n} - D_{1n} = \sigma_{f}$ (normal $\vect{D}$ jumps by the
    free surface charge density).
  \item $B_{2n} - B_{1n} = 0$ (normal $\vect{B}$ is continuous).
  \item $\vect{E}_{2t} - \vect{E}_{1t} = \vect{0}$ (tangential
    $\vect{E}$ is continuous).
  \item $\unitvec{n}\times(\vect{H}_{2} - \vect{H}_{1}) = \vect{K}_{f}$
    (tangential $\vect{H}$ jumps by the free surface current
    density).
\end{itemize}
For a perfect conductor, $\vect{E}$ and $\vect{B}$ inside are zero in
steady state, so at the surface $\vect{E}$ is normal to the surface
and $\vect{B}$ is tangential. The four boundary conditions are the
interface archetype: every reflection, refraction, waveguide mode, and
transmission-line discontinuity in Chapter 10 is solved by enforcing
them.

\begin{mastery}
A conductor is not a perfect conductor at high frequency. The
constitutive law $\vect{J} = \sigma\vect{E}$ combined with the
Amp\`ere-Maxwell and Faraday laws gives, inside a good conductor where
displacement current is negligible, the diffusion equation
$\laplacian\vect{E} = \mu\sigma\,\partial\vect{E}/\partial t$. A
sinusoidal field at angular frequency $\omega$ decays into the conductor
as $\exp(-x/\delta)$ with the \engterm{skin depth}
\[
\delta = \sqrt{\frac{2}{\mu\sigma\omega}}.
\]
For copper at $50\,\si{\hertz}$ the skin depth is about $9\,\si{\milli
\meter}$; at $1\,\si{\mega\hertz}$ it is $65\,\si{\micro\meter}$; at
$1\,\si{\giga\hertz}$ it is about $2\,\si{\micro\meter}$. Current at
radio frequency rides on the surface, so the usable cross-section of a
wire shrinks with frequency, its effective resistance rises, and the
designer plates conductors with thin layers of high-conductivity metal
rather than paying for solid bulk. The same skin depth sets the
thickness of electromagnetic shielding: a shield works once it is
several skin depths thick at the frequency it must block.
\end{mastery}

\section{Worked examples: solenoid, parallel-plate capacitor, coaxial cable}\pathtag{standard}

A long solenoid with $n$ turns per unit length, radius $a$, carrying
current $I$, has an interior field $\vect{B} = \mu_{0} n I\unitvec{z}$
and a negligible exterior field. Its self-inductance per unit length
is computed from the flux linkage per ampere:
\[
\Phi_{\text{per turn}} = B\pi a^{2} = \mu_{0} n I \pi a^{2},
\]
and the linkage per unit length is $n$ turns per metre times
$\Phi_{\text{per turn}}$:
\[
L' = \frac{N\Phi_{\text{per turn}}}{I\ell} = \mu_{0} n^{2} \pi a^{2}.
\]
A solenoid with $n = 1000\,\si{\per\meter}$ and $a = 1\,\si{\centi
\meter}$ has $L' = 4\pi\times 10^{-7}\times 10^{6}\times \pi\times
10^{-4} \approx 0.4\,\si{\milli\henry\per\meter}$. The energy stored
per unit length is $\tfrac{1}{2}L' I^{2}$; at $I = 10\,\si{\ampere}$,
$U/\ell \approx 20\,\si{\milli\joule\per\meter}$.

% Figure: parallel-plate capacitor. Uniform field between plates, charge
% on plates, fringing at the edges.
\begin{figure}[htbp]
\centering
\begin{tikzpicture}[>=Latex, scale=1.0]
% plates
\draw[engred, line width=2pt] (-2.2,1.1) -- (2.2,1.1);
\draw[engblue, line width=2pt] (-2.2,-1.1) -- (2.2,-1.1);
\node[engred, font=\small] at (2.7,1.1) {$+Q$};
\node[engblue, font=\small] at (2.7,-1.1) {$-Q$};
% uniform interior field lines
\foreach \x in {-1.8,-1.2,...,1.9}{
  \draw[enggray, ->] (\x,1.05) -- (\x,-1.05);
}
% fringing at right edge
\draw[enggray] (2.2,0.9) .. controls (2.9,0.6) and (2.9,-0.6) .. (2.2,-0.9);
\draw[enggray] (2.2,0.5) .. controls (3.3,0.2) and (3.3,-0.2) .. (2.2,-0.5);
% fringing at left edge
\draw[enggray] (-2.2,0.9) .. controls (-2.9,0.6) and (-2.9,-0.6) .. (-2.2,-0.9);
\draw[enggray] (-2.2,0.5) .. controls (-3.3,0.2) and (-3.3,-0.2) .. (-2.2,-0.5);
% spacing dimension
\draw[<->] (-2.5,1.1) -- (-2.5,-1.1) node[midway, left, font=\small] {$d$};
% area label
\node[font=\small] at (0,1.45) {plate area $A$};
\node[font=\small] at (0,0) {$E=V/d$};
\node[font=\small] at (0,-1.55) {$C=\varepsilon A/d$};
\end{tikzpicture}
\caption[Parallel-plate capacitor]{The parallel-plate capacitor. Equal
and opposite charges $\pm Q$ on plates of area $A$ at spacing $d$ set up
a uniform interior field $E = V/d$, with the field zero outside in the
ideal limit. A pillbox argument from Gauss's law gives surface charge
$\sigma = \varepsilon E$, so $Q = \varepsilon V A/d$ and the capacitance
is $C = \varepsilon A/d$. The curved fringing field at the edges is the
correction the infinite-plate formula omits; it grows in relative
importance as the gap $d$ approaches the plate dimension.}
\label{fig:vol04ch08:parallel-plate}
\end{figure}


A parallel-plate capacitor with plate area $A$, separation $d$, and
dielectric permittivity $\varepsilon$ has a uniform interior field
$E = V/d$ when held at voltage $V$. The charge on a plate, by Gauss's
law on a pillbox spanning the plate, is $Q = \varepsilon E A =
\varepsilon V A/d$, giving the capacitance
\[
C = \frac{Q}{V} = \frac{\varepsilon A}{d}.
\]
A $1\,\si{\centi\meter\squared}$ plate with $1\,\si{\milli\meter}$
air gap has $C = 8.85\times 10^{-13}\,\si{\farad} = 0.885\,\si{\pico
\farad}$; the picofarad is the natural scale of bench-built
capacitors, and only with thin dielectrics ($d \ll 1\,\si{\milli
\meter}$) and high $\varepsilon_{r}$ (ceramics, electrolytes) do
practical capacitances of microfarads and millifarads become
achievable. The dielectric constants that make these high capacitances
practical are collected in \repopath{data/dielectric\_materials.csv},
which lists relative permittivity and dielectric strength from vacuum
through air, polyethylene, and the class-2 ceramics whose
$\varepsilon_r$ reaches several thousand. The idealised formula
$C = \varepsilon A/d$ omits the fringing field at the plate edges; the
finite-difference solver \repopath{code/laplace\_capacitor.py} relaxes
Laplace's equation on a grid for finite plates and reports the fringing
correction as a function of plate aspect ratio, the quantity the
simulation exercise asks the reader to characterise.

% Figure: coaxial-cable cross section. Inner conductor radius a, outer
% radius b, dielectric between, radial E field and azimuthal B field.
\begin{figure}[htbp]
\centering
\begin{tikzpicture}[>=Latex, scale=1.0]
% dielectric annulus
\fill[engamber!15] (0,0) circle [radius=2.0];
\fill[white] (0,0) circle [radius=0.5];
% conductors
\filldraw[engred] (0,0) circle [radius=0.5];
\draw[engblue, line width=2pt] (0,0) circle [radius=2.0];
% radii labels
\draw[enggray, ->] (0,0) -- (210:0.5) node[midway, below, font=\scriptsize] {$a$};
\draw[enggray, ->] (0,0) -- (150:2.0) node[pos=0.7, above, font=\scriptsize] {$b$};
% radial E field arrows (inward, since inner is +)
\foreach \a in {0,45,...,315}{
  \draw[engred, ->] (\a:1.85) -- (\a:0.65);
}
% azimuthal B field (a dashed circle with arrow)
\draw[engblue, dashed, ->] (30:1.25) arc [start angle=30, end angle=110, radius=1.25];
\node[engblue, font=\scriptsize] at (70:1.5) {$\vect{B}$};
\node[engred, font=\scriptsize] at (-25:1.3) {$\vect{E}$};
\node[white, font=\scriptsize\bfseries] at (0,0) {$+\lambda$};
\node[font=\scriptsize] at (0,-2.5) {dielectric $\varepsilon$, $a<r<b$};
\end{tikzpicture}
\caption[Coaxial-cable cross section]{Cross section of a coaxial cable.
The inner conductor of radius $a$ carries linear charge $+\lambda$ and
current $I$; the outer conductor of inner radius $b$ carries $-\lambda$
and the return current. Between them the electric field is radial,
$E = \lambda/(2\pi\varepsilon r)$, and the magnetic field is azimuthal,
$B = \mu I/(2\pi r)$. Integrating gives capacitance per length
$C' = 2\pi\varepsilon/\ln(b/a)$ and inductance per length
$L' = (\mu/2\pi)\ln(b/a)$, whose ratio fixes the characteristic
impedance $Z_0 = \sqrt{L'/C'}$. The Poynting vector $\vect{E}\times
\vect{B}/\mu$ points along the axis: the power travels in the dielectric
between the conductors, not in the metal.}
\label{fig:vol04ch08:coax-cross-section}
\end{figure}


A coaxial cable of inner radius $a$, outer radius $b$, with the
dielectric of permittivity $\varepsilon$ between, carries a linear
charge $\lambda$ on the inner conductor (and $-\lambda$ on the outer).
Gauss's law on a cylindrical Gaussian surface gives
\[
E(r) = \frac{\lambda}{2\pi \varepsilon r}, \quad a<r<b.
\]
The potential difference between the conductors is
\[
V = -\int_{b}^{a} E\,\dd r = \frac{\lambda}{2\pi \varepsilon}\ln(b/a),
\]
and the capacitance per unit length is
\[
C' = \frac{\lambda}{V} = \frac{2\pi \varepsilon}{\ln(b/a)}.
\]
A polyethylene-dielectric cable with $\varepsilon_{r} = 2.25$, $b/a =
3$, has $C' = 2\pi\varepsilon_{0}\times 2.25/\ln 3 \approx 114\,\si{
\pico\farad\per\meter}$, the standard $50\,\si{\ohm}$ RG-58 cable.
The inductance per unit length, computed from Amp\`ere's law and the
flux linkage between the conductors, is
\[
L' = \frac{\mu}{2\pi}\ln(b/a),
\]
giving a characteristic impedance
\[
Z_{0} = \sqrt{L'/C'} = \frac{1}{2\pi}\sqrt{\mu/\varepsilon}\,\ln(b/a).
\]
For air dielectric, $Z_{0} \approx 60\,\ln(b/a)$; for polyethylene,
$Z_{0} \approx 40\,\ln(b/a)$, recovering $50\,\si{\ohm}$ at $b/a
\approx 3$. The characteristic impedance is the load that, when
attached at the cable's far end, absorbs all incident power; mismatched
loads produce reflections, treated in Chapter 10. The script
\repopath{code/coax\_parameters.py} evaluates $C'$, $L'$, $Z_0$, and the
propagation velocity for any geometry and dielectric; for the RG-58
parameters it returns $Z_0 \approx 47\,\si{\ohm}$ and a propagation
velocity of $0.67c$, the velocity factor printed on the cable's
datasheet.

\section{Failure: ground loops, induced EMI, lightning}\pathtag{core}

Every electrical system in the world is embedded in an
electromagnetic environment that includes power-line magnetic fields,
radio-frequency signals, and atmospheric electrical activity. Three
failure mechanisms recur across every domain in which fields are
measured or signals are transmitted, and a chapter that does not name
them has not described real electromagnetism.

A \engterm{ground loop} arises when two devices that share a signal
connection also share separate paths to earth. The two earth paths,
typically through separate building grounds, are at slightly different
potentials because of currents flowing through the building's
distribution system. The voltage difference drives a current through
the signal connection between the devices, contaminating any
small-signal measurement that rides on top of it. The mechanism is
Faraday's law applied to the area enclosed by the two ground paths
and the signal cable: time-varying magnetic flux through that area
induces an EMF around the loop. Mains-frequency hum at $50\,\si{\hertz}$
or $60\,\si{\hertz}$ in audio recordings, oscilloscope traces, and
biological signal recordings is the everyday signature.

The cure is a single-point ground: every device's earth connection is
brought to a single physical node, so no loop area exists. When a
single-point ground cannot be achieved (long signal runs, multiple
buildings), the working tools are differential signalling (the signal
is the difference between two conductors, and common-mode interference
on both is rejected), optical isolation (the signal is converted to
light and back, breaking the electrical connection), and transformer
coupling (the signal is transferred through a magnetic core, again
breaking the direct electrical connection).

\engterm{Electromagnetic interference} (EMI) is unintended coupling of
electromagnetic energy from a source into a victim circuit. The source
may be a switching power supply whose high-$dI/dt$ current pulses
radiate; a digital clock line whose square-wave edges have harmonic
content into the gigahertz; a motor whose commutation produces
broadband sparks; a radio transmitter operating nearby. The coupling
mechanism may be capacitive (high-impedance victim sees high-$dV/dt$
source through stray capacitance), inductive (low-impedance victim
sees high-$dI/dt$ source through mutual inductance), or radiative
(the source is far enough away that near-field coupling has fallen
off and the wave propagates as a plane wave; Chapter 10). The
frequency band of the source sets which coupling dominates and which
wavelength matters; the table \repopath{data/em\_spectrum\_bands.csv}
lists the engineering bands from power-line and audio frequencies
through the radio spectrum to visible light, with the wavelength range
of each, so that a conductor's length can be compared against the
wavelength of the interference it carries.

Two principles govern EMI mitigation. First, the source is reduced at
the source: a switching converter's $dI/dt$ is shaped by snubber
networks, slew-rate limits, or spread-spectrum modulation that
spreads the harmonic content over a band rather than concentrating it
at a single frequency. Second, the victim is shielded and the loop
areas of victim circuits are minimised: a twisted pair carrying the
signal and its return places the two conductors close together, so
the loop area through which external flux passes is small, and any
flux that does couple in produces equal EMF in both conductors that
cancels at a differential receiver.

A representative regulation is the European Union's EMC Directive
(2014/30/EU) and the corresponding harmonised standards in the
EN 61000 and EN 55000 series, which specify radiated and conducted
emission limits, immunity test levels, and the certification
procedure for placing electrical equipment on the European market.
Equivalent regulation in the United States is FCC Part 15 for
unintentional radiators. Compliance is a working constraint on the
design of consumer electronics, industrial equipment, and medical
devices.

\engterm{Lightning} is the failure mode at the high end of natural
electrical activity. A typical cloud-to-ground stroke carries
$30\,\si{\kilo\ampere}$ peak current over a microsecond rise time,
producing $dI/dt \sim 10^{10}\,\si{\ampere\per\second}$ in the
channel and very large radiated fields at distances of hundreds of
metres. The direct strike to a structure is the rare event; the
indirect coupling, in which the magnetic field of the stroke induces
EMF in every loop in nearby equipment, is the common one.
Engineering protection against lightning is a layered discipline:
\begin{itemize}[leftmargin=*]
  \item A \engterm{lightning protection system} provides a preferred
    path for the strike current from the structure's roof to a
    grounding electrode buried at depth, avoiding the interior of the
    structure. International standards IEC 62305 specify the design.
  \item \engterm{Surge protective devices} (SPDs) clamp the voltage on
    incoming power, signal, and antenna lines to a safe value during
    the transient, diverting the surge current to ground. Three SPD
    classes are defined by IEC 61643: Type 1 for direct-strike currents
    at the service entrance, Type 2 for indirect surges in the
    distribution panel, Type 3 for sensitive equipment at the outlet.
  \item \engterm{Bonding} of every grounded structure (panel, conduit,
    rack, equipment chassis) to a common reference at the service
    entrance ensures that during a strike, the entire structure
    floats together rather than allowing potential differences that
    would drive currents through interior loops.
\end{itemize}

A common failure mode of grid-tied photovoltaic installations is
surge damage from indirect lightning: induced EMF in the long DC
cabling between modules and inverter drives currents that exceed the
inverter input ratings, even when no direct strike occurred. The cure
is Type 2 SPDs at the inverter input, properly bonded to the same
grounding system as the array's frame. The same mechanism applies to
wind turbines, where the long blade acts as an antenna for nearby
lightning, and to telecommunications towers, where lightning is a
routine operational hazard rather than an exceptional event.

\begin{principle}[Every cable is an antenna]
Every conductor in an electrical system both radiates and receives
electromagnetic energy at frequencies whose wavelength is comparable
to or larger than the conductor's length. The engineer who treats a
wire as merely a transport for current is correct only when the
wire is short compared to a wavelength; outside that regime, the wire
is an antenna whether the designer intended it to be one or not. The
discipline of EMC, lightning protection, and signal integrity is the
discipline of taking unintended antennas seriously.
\end{principle}

\section{Project}\pathtag{core}

\begin{project}[Magnetic field around a current-carrying wire]
Build a bench measurement of the magnetic field around a straight
current-carrying wire and verify the $1/r$ law predicted by
Amp\`ere's law.

\textbf{Materials.}
A bench DC power supply (current-limited, $0$ to $5\,\si{\ampere}$ at
$\le 10\,\si{\volt}$); a $1\,\si{\meter}$ straight length of $14$-gauge
copper wire mounted vertically; a Hall-effect magnetometer (a
calibrated probe with a sensitivity of $\sim 10\,\si{\micro\tesla}$ or
better, available as a USB-connected instrument for under \$200, or a
smartphone with a built-in magnetometer and a free measurement app); a
metre rule.

\textbf{Procedure.}
\begin{enumerate}
  \item Calibrate the magnetometer to the local geomagnetic field by
    pointing it in three orthogonal directions and recording the
    components. The Earth's field magnitude at the bench is the
    quadratic sum; verify against a published model (NOAA WMM,
    current as of 2025) within $5\,\%$.
  \item Set the supply to $1.0\,\si{\ampere}$ and measure the
    magnetic-field component perpendicular to the wire at radial
    distances $r = 1, 2, 3, 5, 10, 20, 50\,\si{\centi\meter}$ from
    the wire's centre. Subtract the geomagnetic background.
  \item Repeat at $2.0\,\si{\ampere}$ and $3.0\,\si{\ampere}$.
  \item Plot $B(r)$ on log-log axes; fit a power law $B = A r^{-n}$
    and report $n$ and $A$. Compare to the predicted $\mu_{0}I/(2\pi
    r)$; report the relative error.
\end{enumerate}

\textbf{Report.} The deliverable is a one-page report with the
calibration data, the measured $B(r)$ at each current, the fit
parameters, and a discussion of the dominant error sources. Expected
sources: geomagnetic-background instability ($\sim 1\,\si{\micro
\tesla}$ during the measurement window), probe-position uncertainty
(largest at small $r$), supply-current uncertainty
($\sim 1\,\%$), and finite-wire-length corrections (at $r$ comparable
to or larger than the wire's length, the field falls faster than
$1/r$). The $1/r$ law should be recovered with $n = 1.00\pm 0.05$
across the central two decades of $r$. Failure to recover $n = 1$ at
either end of the range is the project's pedagogic outcome.
\end{project}

\section{Exercises}

\subsection*{Calculation}\pathtag{core}

\begin{exercise}
Two point charges of $+2\,\si{\micro\coulomb}$ and $-3\,\si{\micro
\coulomb}$ are separated by $30\,\si{\centi\meter}$ in air. Compute
the electric-field magnitude at the midpoint between them, and the
force one would experience by a $1\,\si{\micro\coulomb}$ test charge
placed at the midpoint.

\paragraph{Solution sketch.} The midpoint is $0.15\,\si{\meter}$ from
each charge. Each field magnitude is $E = q/(4\pi\varepsilon_0 r^2) =
(8.99\times 10^9)(q)/(0.15)^2$: the $+2\,\si{\micro\coulomb}$ gives
$E_1 = 7.99\times 10^{5}\,\si{\volt\per\meter}$ pointing away from the
positive charge, the $-3\,\si{\micro\coulomb}$ gives $E_2 = 1.20\times
10^{6}\,\si{\volt\per\meter}$ pointing toward the negative charge. Both
point in the same direction along the line (away from $+$, toward $-$),
so they add: $E = 2.0\times 10^{6}\,\si{\volt\per\meter}$. The force on
the test charge is $F = qE = (10^{-6})(2.0\times 10^{6}) = 2.0\,\si{
\newton}$, directed from the positive toward the negative charge. The
sign trap is to subtract rather than add: along the joining line the
two contributions reinforce.
\end{exercise}

\begin{exercise}
A parallel-plate capacitor has plates of area $50\,\si{\centi\meter
\squared}$ separated by $0.5\,\si{\milli\meter}$ of air. Compute the
capacitance. The plates are charged to $200\,\si{\volt}$; compute the
charge on each plate, the energy stored, and the force of attraction
between the plates.

\paragraph{Solution sketch.} With $A = 50\times 10^{-4}\,\si{\meter
\squared}$ and $d = 5\times 10^{-4}\,\si{\meter}$, $C = \varepsilon_0
A/d = (8.85\times 10^{-12})(5\times 10^{-3})/(5\times 10^{-4}) =
8.85\times 10^{-11}\,\si{\farad} = 88.5\,\si{\pico\farad}$. The charge
is $Q = CV = (8.85\times 10^{-11})(200) = 1.77\times 10^{-8}\,\si{
\coulomb} = 17.7\,\si{\nano\coulomb}$. The stored energy is $U =
\tfrac12 CV^2 = \tfrac12(8.85\times 10^{-11})(200)^2 = 1.77\times
10^{-6}\,\si{\joule}$. The attractive force is $F = \tfrac12
\varepsilon_0 E^2 A$ with $E = V/d = 4\times 10^{5}\,\si{\volt\per
\meter}$, giving $F = \tfrac12(8.85\times 10^{-12})(4\times 10^5)^2
(5\times 10^{-3}) = 3.5\times 10^{-3}\,\si{\newton}$, equivalently
$F = U/d = QE/2$. The factor of one half in the force is the common
error: the field acting on one plate is the field of the \emph{other}
plate, $\sigma/(2\varepsilon_0)$, not the full interior field.
\end{exercise}

\begin{exercise}
A solenoid is wound with $500$ turns over a length of $25\,\si{\centi
\meter}$ on a $2\,\si{\centi\meter}$-radius non-magnetic core. With
$3\,\si{\ampere}$ flowing, compute the interior magnetic field and the
self-inductance.

\paragraph{Solution sketch.} The turn density is $n = 500/0.25 = 2000\,
\si{\per\meter}$. The interior field is $B = \mu_0 n I = (4\pi\times
10^{-7})(2000)(3) = 7.54\times 10^{-3}\,\si{\tesla} = 7.5\,\si{\milli
\tesla}$. The inductance follows from flux linkage: $L = \mu_0 n^2 A
\ell$ with $A = \pi(0.02)^2 = 1.26\times 10^{-3}\,\si{\meter\squared}$
and $\ell = 0.25\,\si{\meter}$, giving $L = (4\pi\times 10^{-7})(2000)^2
(1.26\times 10^{-3})(0.25) = 1.58\times 10^{-3}\,\si{\henry} \approx
1.6\,\si{\milli\henry}$. Equivalently $L = N\Phi/I = N B A/I$ with
$N = 500$. The check is that $\tfrac12 LI^2 = 7.1\,\si{\milli\joule}$
matches $\tfrac12(B^2/\mu_0)\times(\text{interior volume})$.
\end{exercise}

\begin{exercise}
A coaxial cable has inner conductor of radius $0.4\,\si{\milli\meter}$
and outer conductor of inner radius $2.0\,\si{\milli\meter}$, with
polyethylene dielectric ($\varepsilon_{r} = 2.25$, $\mu_{r} = 1$).
Compute the capacitance per metre, the inductance per metre, and the
characteristic impedance.

\paragraph{Solution sketch.} The ratio is $b/a = 2.0/0.4 = 5$, so
$\ln(b/a) = 1.609$. Capacitance per metre: $C' = 2\pi\varepsilon_0
\varepsilon_r/\ln(b/a) = 2\pi(8.85\times 10^{-12})(2.25)/1.609 =
7.78\times 10^{-11}\,\si{\farad\per\meter} = 77.8\,\si{\pico\farad\per
\meter}$. Inductance per metre: $L' = (\mu_0/2\pi)\ln(b/a) = (2\times
10^{-7})(1.609) = 3.22\times 10^{-7}\,\si{\henry\per\meter} = 322\,\si{
\nano\henry\per\meter}$. Characteristic impedance: $Z_0 = \sqrt{L'/C'}
= \sqrt{3.22\times 10^{-7}/7.78\times 10^{-11}} = 64\,\si{\ohm}$;
equivalently $Z_0 = (60/\sqrt{\varepsilon_r})\ln(b/a) = (60/1.5)(1.609)
= 64\,\si{\ohm}$. The script \repopath{code/coax\_parameters.py}
reproduces these values.
\end{exercise}

\begin{exercise}
A circular loop of radius $5\,\si{\centi\meter}$ carries a current that
ramps linearly from $0$ to $2\,\si{\ampere}$ in $10\,\si{\milli\second}$.
Compute the rate of change of the magnetic flux through a concentric
loop of radius $1\,\si{\centi\meter}$, and the EMF induced around the
inner loop (assume the inner loop is in the central region where the
outer loop's field is approximately uniform and given by $B =
\mu_{0}I/(2R)$).

\paragraph{Solution sketch.} The outer loop's centre field is $B =
\mu_0 I/(2R)$ with $R = 0.05\,\si{\meter}$. The flux through the inner
loop of radius $r = 0.01\,\si{\meter}$, area $a = \pi r^2 = 3.14\times
10^{-4}\,\si{\meter\squared}$, is $\Phi = Ba = \mu_0 I a/(2R)$. Its
rate of change tracks $\dd I/\dd t = 2/(10\times 10^{-3}) = 200\,\si{
\ampere\per\second}$: $\dd\Phi/\dd t = \mu_0 a/(2R)\cdot\dd I/\dd t =
(4\pi\times 10^{-7})(3.14\times 10^{-4})/(0.1)\times 200 = 7.9\times
10^{-7}\,\si{\weber\per\second}$. The induced EMF around the single
inner loop is $\mathcal{E} = -\dd\Phi/\dd t = -0.79\,\si{\micro\volt}$,
a magnitude of about $0.8\,\si{\micro\volt}$. The smallness of the
number is the lesson: tightly coupled single turns at modest $\dd I/
\dd t$ induce microvolts, which is why transformers use many turns and
high frequencies.
\end{exercise}

\begin{exercise}
The solar irradiance at Earth's surface on a clear day is around
$1000\,\si{\watt\per\meter\squared}$. Compute the peak electric-field
amplitude and the peak magnetic-field amplitude of the plane-wave
component.

\paragraph{Solution sketch.} From $\langle S\rangle = \tfrac12
\varepsilon_0 c E_0^2$, solve for the peak electric field: $E_0 =
\sqrt{2\langle S\rangle/(\varepsilon_0 c)} = \sqrt{2(1000)/((8.85\times
10^{-12})(3\times 10^8))} = 868\,\si{\volt\per\meter}$. The peak
magnetic field is $B_0 = E_0/c = 868/(3\times 10^8) = 2.9\times
10^{-6}\,\si{\tesla} = 2.9\,\si{\micro\tesla}$, comparable in magnitude
to the Earth's static field but oscillating at optical frequency. The
script \repopath{code/plane\_wave.py} returns the same values and the
$1361\,\si{\watt\per\meter\squared}$ space-side case.
\end{exercise}

\subsection*{Derivation}\pathtag{standard}

\begin{exercise}
Starting from $\oint \vect{E}\cdot \dd\vect{A} = Q_{\text{enc}}/
\varepsilon_{0}$, derive Coulomb's law for the field of a point
charge. State the symmetry assumption that closes the integral.

\paragraph{Solution sketch.} The symmetry assumption is that a point
charge at the origin produces a field that is purely radial and depends
only on $r$, so that on a sphere of radius $r$ centred on the charge,
$\vect{E}\cdot\dd\vect{A} = E(r)\,\dd A$ with $E(r)$ constant over the
sphere. The flux integral then factors: $\oint E(r)\,\dd A = E(r)\cdot
4\pi r^2$. Setting this equal to $q/\varepsilon_0$ gives $E(r) =
q/(4\pi\varepsilon_0 r^2)$, and $\vect{F} = q'\vect{E}$ on a second
charge recovers Coulomb's law $F = q q'/(4\pi\varepsilon_0 r^2)$. The
derivation uses no integration beyond the area of a sphere; the entire
content is the symmetry argument that makes $E$ constant on the chosen
surface.
\end{exercise}

\begin{exercise}
Derive the displacement current term in the Amp\`ere-Maxwell law by
demanding consistency between Amp\`ere's law without displacement
current and the continuity equation $\partial \rho/\partial t +
\divg \vect{J} = 0$. State precisely where the bare Amp\`ere law
fails.

\paragraph{Solution sketch.} The bare law is $\curl\vect{B} = \mu_0
\vect{J}$. Take the divergence of both sides: $\divg(\curl\vect{B}) =
\mu_0\divg\vect{J}$. The left side is identically zero, since the
divergence of any curl vanishes. So the bare law forces $\divg\vect{J}
= 0$ everywhere and always. Continuity says $\divg\vect{J} = -\partial
\rho/\partial t$, which is non-zero wherever charge accumulates, as on a
charging capacitor plate. The contradiction is exactly there: the bare
law fails whenever $\partial\rho/\partial t \neq 0$. Adding a term
$\mu_0\varepsilon_0\,\partial\vect{E}/\partial t$ and using Gauss's law
$\divg\vect{E} = \rho/\varepsilon_0$ gives $\divg(\mu_0\vect{J} +
\mu_0\varepsilon_0\,\partial\vect{E}/\partial t) = \mu_0(-\partial\rho/
\partial t) + \mu_0\,\partial\rho/\partial t = 0$ identically. The
divergence is now consistent, and the added term is the displacement
current.
\end{exercise}

\begin{exercise}
Derive the Poynting theorem $\partial (u_{e}+u_{m})/\partial t +
\divg \vect{S} = -\vect{J}\cdot\vect{E}$ by combining $\vect{E}\cdot
(\curl \vect{B} = \mu_{0}\vect{J} + \mu_{0}\varepsilon_{0}\partial
\vect{E}/\partial t)$ with $\vect{B}\cdot(\curl \vect{E} = -\partial
\vect{B}/\partial t)$ and using the vector identity $\divg(\vect{E}
\times \vect{B}) = \vect{B}\cdot \curl \vect{E} - \vect{E}\cdot \curl
\vect{B}$.

\paragraph{Solution sketch.} Dot the Amp\`ere-Maxwell law with $\vect{E}$:
$\vect{E}\cdot\curl\vect{B} = \mu_0\vect{J}\cdot\vect{E} + \mu_0
\varepsilon_0\vect{E}\cdot\partial\vect{E}/\partial t$. Dot Faraday's
law with $\vect{B}$: $\vect{B}\cdot\curl\vect{E} = -\vect{B}\cdot
\partial\vect{B}/\partial t$. Subtract the first from the second and
apply the identity, $\divg(\vect{E}\times\vect{B}) = \vect{B}\cdot
\curl\vect{E} - \vect{E}\cdot\curl\vect{B}$, on the left. Recognise the
two time-derivative products as $\vect{E}\cdot\partial\vect{E}/\partial
t = \tfrac12\partial|\vect{E}|^2/\partial t$ and likewise for $\vect{B}$.
Divide through by $\mu_0$ and collect: $\partial/\partial t(\tfrac12
\varepsilon_0|\vect{E}|^2 + \tfrac{1}{2\mu_0}|\vect{B}|^2) + \divg(
\vect{E}\times\vect{B}/\mu_0) = -\vect{J}\cdot\vect{E}$, which is the
Poynting theorem with $\vect{S} = \vect{E}\times\vect{B}/\mu_0$. The
right side is the volumetric work the field does on charge, the resistive
loss $\vect{J}\cdot\vect{E}$.
\end{exercise}

\begin{exercise}
Derive the four boundary conditions at the interface between two
linear isotropic media by applying Maxwell's equations in integral
form to a pillbox (for normal components) and to a rectangular loop
(for tangential components) straddling the interface, in the limit
that the pillbox height and the loop's normal dimension go to zero.

\paragraph{Solution sketch.} Normal components: apply $\oint\vect{D}\cdot
\dd\vect{A} = Q_{f,\text{enc}}$ to a flat pillbox of face area $\dd A$
and vanishing height. As the height goes to zero the side-wall flux
vanishes and only the two faces contribute, giving $(D_{2n} - D_{1n})
\,\dd A = \sigma_f\,\dd A$, so $D_{2n} - D_{1n} = \sigma_f$. The same
pillbox with $\divg\vect{B}=0$ gives $B_{2n} - B_{1n} = 0$. Tangential
components: apply $\oint\vect{E}\cdot\dd\vect{\ell} = -\dd\Phi_B/\dd t$
to a thin rectangular loop straddling the interface with long sides
$\dd\ell$ tangential and short sides normal and vanishing. As the loop
area goes to zero the flux term vanishes, leaving $(\vect{E}_{2t} -
\vect{E}_{1t})\,\dd\ell = 0$, so tangential $\vect{E}$ is continuous.
The same loop with the Amp\`ere-Maxwell law retains the enclosed free
surface current $\vect{K}_f$ as the area shrinks, giving $\unitvec{n}
\times(\vect{H}_2 - \vect{H}_1) = \vect{K}_f$. The trick in each case is
that one dimension of the Gaussian region is sent to zero, killing the
volume or area term and leaving only the surface source.
\end{exercise}

\begin{exercise}
Show that the electromagnetic field's momentum density is
$\vect{g} = \varepsilon_{0}\vect{E}\times \vect{B} = \vect{S}/c^{2}$,
and that radiation pressure on a perfectly absorbing surface is
$P = I/c$ where $I$ is the time-averaged intensity. Compute the
radiation pressure of solar irradiance at $1361\,\si{\watt\per\meter
\squared}$.

\paragraph{Solution sketch.} The field momentum density is the Poynting
vector divided by $c^2$: $\vect{g} = \vect{S}/c^2 = \varepsilon_0
\vect{E}\times\vect{B}$, using $\vect{S} = \vect{E}\times\vect{B}/\mu_0$
and $1/(\mu_0 c^2) = \varepsilon_0$. The momentum delivered per unit
time per unit area to an absorbing surface is the momentum flux, which
equals the energy flux divided by $c$ for light, so the pressure is
$P = \langle S\rangle/c = I/c$. (A perfect reflector doubles this.)
Numerically $P = 1361/(3\times 10^8) = 4.5\times 10^{-6}\,\si{\pascal}$,
about $4.5\,\si{\micro\pascal}$: tiny on a hand, but the basis of solar
sails and a real perturbation on spacecraft attitude over long
missions. The script \repopath{code/plane\_wave.py} reports the same
pressure.
\end{exercise}

\subsection*{Estimation}\pathtag{core}

\begin{exercise}
Estimate the energy stored in the magnetic field of the Earth in the
volume from the surface to $1$ Earth radius altitude, treating the
field as roughly dipolar with surface magnitude $50\,\si{\micro\tesla}$.
State assumptions about the field's variation with altitude.

\paragraph{Solution sketch.} Energy density is $u_m = B^2/(2\mu_0)$. At the
surface, $u_m = (50\times 10^{-6})^2/(2\times 4\pi\times 10^{-7}) =
1.0\times 10^{-3}\,\si{\joule\per\meter\cubed}$. A dipole field falls as
$1/r^3$, so $u_m$ falls as $1/r^6$; most of the energy in the shell sits
near the surface. Take the shell volume from $R_E = 6.4\times 10^6\,\si{
\meter}$ to $2R_E$ as roughly $V \sim 4\pi R_E^2\cdot R_E = 4\pi R_E^3
\approx 3\times 10^{21}\,\si{\meter\cubed}$, but weight by the $1/r^6$
falloff, which cuts the effective volume to a fraction of order $0.3$ of
a one-radius shell. An order-of-magnitude estimate is $U \sim u_m\cdot
0.3\cdot 4\pi R_E^3/3 \sim (10^{-3})(0.3)(1.1\times 10^{21}) \sim
3\times 10^{17}\,\si{\joule}$. The dominant assumption is the dipole
$1/r^3$ law, which the real magnetosphere violates above a few Earth
radii where the solar wind distorts it; the estimate is good to a factor
of a few, no better.
\end{exercise}

\begin{exercise}
Estimate the peak voltage induced across the terminals of a
$1\,\si{\meter\squared}$ open loop placed $100\,\si{\meter}$ from a
lightning channel carrying a peak current of $30\,\si{\kilo\ampere}$
with rise time $1\,\si{\micro\second}$. State the dominant
uncertainty in the answer.

\paragraph{Solution sketch.} The channel field at $r = 100\,\si{\meter}$
is $B = \mu_0 I/(2\pi r) = (4\pi\times 10^{-7})(30\times 10^3)/(2\pi
\times 100) = 6.0\times 10^{-5}\,\si{\tesla}$ at peak. The field rises in
$\tau = 1\,\si{\micro\second}$, so $\dd B/\dd t \approx B/\tau = 6.0
\times 10^{-5}/10^{-6} = 60\,\si{\tesla\per\second}$. The EMF in a
$1\,\si{\meter\squared}$ loop is $\mathcal{E} = A\,\dd B/\dd t = 60\,\si{
\volt}$, of order tens of volts. The dominant uncertainty is the loop's
orientation and effective area relative to the field, which can swing the
flux linkage by a full order of magnitude depending on geometry; the
$1/r$ near-field approximation and the assumed linear current ramp add
smaller errors. The lesson is that a modest loop a hundred metres from a
strike sees tens of volts with no direct contact, which is why induced
surges, not direct strikes, are the common failure.
\end{exercise}

\begin{exercise}
Estimate the time it would take for the field of a $1\,\si{\farad}$
capacitor charged to $100\,\si{\volt}$ to discharge through the
leakage resistance of typical air at $50\,\%$ relative humidity. State
the assumptions and report a factor-of-ten confidence range.

\paragraph{Solution sketch.} The discharge time constant is $\tau = RC$.
The capacitance is $C = 1\,\si{\farad}$. The leakage resistance depends
on the geometry and the path through air; take a representative bench
gap of $1\,\si{\centi\meter}$ over $1\,\si{\centi\meter\squared}$ of
plate and air conductivity in the range $\sigma \sim 10^{-14}$ to
$10^{-13}\,\si{\siemens\per\meter}$ at moderate humidity. Then $R =
\ell/(\sigma A) = 10^{-2}/(10^{-13}\times 10^{-4}) \sim 10^{15}\,\si{
\ohm}$. The time constant is $\tau = RC \sim 10^{15}\,\si{\second}$,
which is enormous, longer than the age of the universe, so air leakage
across a clean gap is not what discharges a real capacitor. The
assumption that fails is that air across the gap is the path; in
practice surface leakage across the dielectric, the dielectric's own
bulk conductivity, and contamination dominate, dropping the practical
time constant to days or hours. The factor-of-ten range on the air-only
estimate spans $10^{14}$ to $10^{16}\,\si{\second}$, and the point of
the exercise is that the idealised number is absurd: the real leakage
path is never clean air.
\end{exercise}

\begin{exercise}
Estimate the radio-frequency energy radiated by a $100\,\si{\watt}$
incandescent bulb at the bulb's filament-fluctuation frequencies.
State the bandwidth assumed and the radiation efficiency at those
frequencies.

\paragraph{Solution sketch.} An incandescent filament runs near constant
temperature; its electrical fluctuations are dominated by the
power-line frequency and its low harmonics ($50$ or $60\,\si{\hertz}$
and a few multiples) plus thermal noise. Take a fluctuation bandwidth
of order $1\,\si{\kilo\hertz}$ and a fluctuating power that is a small
fraction, say $1\,\%$, of the $100\,\si{\watt}$, so $\sim 1\,\si{\watt}$
of fluctuating electrical power. The filament is electrically tiny
compared to the wavelength at these frequencies ($\lambda \sim
6000\,\si{\kilo\meter}$ at $50\,\si{\hertz}$), and the radiation
efficiency of an antenna of length $\ell \ll \lambda$ scales as
$(\ell/\lambda)^2$, which is of order $(0.01/6\times 10^6)^2 \sim
10^{-18}$. The radiated RF power is then $\sim 1\,\si{\watt}\times
10^{-18} \sim 10^{-18}\,\si{\watt}$: negligible. The lesson is that
low-frequency sources radiate almost nothing because they are tiny
compared to a wavelength; the EMI problem lives at high frequency,
where the same conductor becomes a fair fraction of a wavelength.
\end{exercise}

\subsection*{Simulation}\pathtag{mastery}

\begin{exercise}
Implement a finite-difference solver for Laplace's equation $\laplacian
\phi = 0$ in two dimensions on a square grid, with Dirichlet boundary
conditions corresponding to a parallel-plate capacitor with finite-
size plates. Compare the computed capacitance against the
infinite-plate formula and identify the fringing-field correction as
a function of plate aspect ratio.

\paragraph{Solution sketch.} Discretise $\phi$ on an $N\times N$ grid and
relax the five-point stencil $\phi_{ij} = \tfrac14(\phi_{i+1,j} +
\phi_{i-1,j} + \phi_{i,j+1} + \phi_{i,j-1})$ at every free node, holding
the two plate rows at $\pm V$ and the box boundary at zero, iterating
until the maximum change per sweep falls below tolerance. Recover the
capacitance from the stored energy, $C = 2U/V^2$ with $U = \tfrac12
\varepsilon_0\sum|\grad\phi|^2\,\dd A$. The computed $C$ exceeds the
infinite-plate $\varepsilon_0 A/d$ because the fringing field stores
extra energy beyond the plate edges; the surplus is a larger fraction
of the total at small aspect ratio (plate width over gap) and shrinks
toward zero as the plates grow wide. The reference implementation is
\repopath{code/laplace\_capacitor.py}, which prints $C/C_{\infty}$ for
several aspect ratios and converges in a few thousand Gauss-Seidel
sweeps. The expected result is a fringing surplus of tens of percent at
aspect ratio of a few, falling as the plates widen.
\end{exercise}

\begin{exercise}
Implement a numerical evaluation of the Biot-Savart integral for a
finite straight wire segment, and compare against the analytical
$B(r,\theta) = (\mu_{0} I/4\pi r)(\sin\theta_{2}-\sin\theta_{1})$
result. Use the simulation to map the field around a square current
loop and identify the centre-of-loop value.

\paragraph{Solution sketch.} Discretise the wire into many short elements
$\dd\vect{\ell}$ and sum $\dd\vect{B} = (\mu_0 I/4\pi)\,\dd\vect{\ell}
\times\unitvec{r}/r^2$ over them with midpoint quadrature. For a single
finite segment seen perpendicular at distance $r$, the numerical sum
must match $(\mu_0 I/4\pi r)(\sin\theta_2 - \sin\theta_1)$, and it
approaches the infinite-wire $\mu_0 I/(2\pi r)$ as the segment
lengthens. For a square loop of side $a$, sum the four sides; the centre
value should converge to $2\sqrt{2}\,\mu_0 I/(\pi a)$, which is four
times the contribution of one side evaluated at the loop centre. The
reference implementation \repopath{code/biot\_savart\_segment.py}
reproduces both checks; the convergence rate is set by the number of
quadrature points and is poorest where the field point is closest to the
wire, since the integrand peaks there.
\end{exercise}

\subsection*{Design}\pathtag{standard}

\begin{exercise}
Design a single-layer air-core solenoid that produces $10\,\si{\milli
\tesla}$ on its axis with no more than $20\,\si{\watt}$ dissipation.
Specify the wire gauge, number of turns, length, radius, and supply
current. State the dominant thermal-management constraint and the
mass of copper required.

\paragraph{Solution sketch.} Target $B = 10\,\si{\milli\tesla} = \mu_0 n
I$, so $nI = B/\mu_0 = 10^{-2}/(4\pi\times 10^{-7}) = 7.96\times 10^3\,
\si{\ampere\per\meter}$ (amp-turns per metre). Pick a workable
operating point: $I = 4\,\si{\ampere}$ needs $n \approx 2000\,\si{\per
\meter}$, which suits a single layer of about $0.5\,\si{\milli\meter}$
enamelled wire (AWG 24, roughly $0.51\,\si{\milli\meter}$ overall).
Choose length $\ell = 0.2\,\si{\meter}$ (so $N = n\ell = 400$ turns) and
radius $a = 2\,\si{\centi\meter}$. The wire length is $L_w = N\cdot 2\pi
a = 400\times 0.126 = 50\,\si{\meter}$; AWG 24 has resistance about
$0.084\,\si{\ohm\per\meter}$, so $R \approx 4.2\,\si{\ohm}$ and the
dissipation is $I^2 R = 16\times 4.2 = 67\,\si{\watt}$, over the
$20\,\si{\watt}$ budget. Bring it under budget by lowering current and
raising turns, or by using thicker wire: dropping to $I = 2\,\si{
\ampere}$ with $n = 4000\,\si{\per\meter}$ gives $I^2 R \approx 17\,\si{
\watt}$ for the same field. The dominant constraint is heat removal from
a long thin winding with poor surface-to-volume ratio; natural
convection limits the surface flux to a few hundred $\si{\watt\per\meter
\squared}$, so the winding surface must exceed roughly $0.1\,\si{\meter
\squared}$ or forced cooling is needed. The copper mass is $\rho_{Cu}
\times L_w\times A_{wire} = 8960\times 50\times (2\times 10^{-7})
\approx 0.09\,\si{\kilogram}$.
\end{exercise}

\begin{exercise}
Design a single-point ground for a small audio-recording studio
containing a microphone preamp, a computer audio interface, and a
monitor amplifier. Sketch the grounding topology, identify potential
ground loops, and specify the bonding cable gauge.

\paragraph{Solution sketch.} Bring every chassis earth to one physical node
(a star point, typically the audio interface or a dedicated bus bar)
with separate radial conductors, so no two devices share a return path
that could enclose loop area. The ground loops to watch are: the preamp
and interface both earthed through their own mains plugs while also
joined by the audio cable shield (cut the shield at one end, or use a
balanced connection); the computer earthed through its own supply and
through USB to the interface; the monitor amp earthed through its mains
plug and through the speaker-level or line return. Use balanced
(differential) interconnects between devices wherever the equipment
supports them, so common-mode hum is rejected at the receiver. The
bonding conductors carry no signal current, only fault and equalising
current, so a modest gauge suffices for the bonds themselves; a common
choice is $2.5\,\si{\milli\meter\squared}$ (about AWG 14) for the star
bonds, sized for fault current and mechanical robustness rather than
for the milliampere signal currents. The deliverable is a star diagram
with one node and radial bonds, plus the note that the safety earth on
each mains plug must not be defeated to break a loop.
\end{exercise}

\begin{exercise}
Design a Type 2 surge protective device specification for a $5\,\si{
\kilo\watt}$ residential photovoltaic inverter's DC input. Specify
the maximum continuous operating voltage, nominal discharge current,
voltage protection level, and response time. Reference IEC 61643-31
or the equivalent national standard.

\paragraph{Solution sketch.} A $5\,\si{\kilo\watt}$ residential array
typically runs at a few hundred volts DC; size the SPD's maximum
continuous operating voltage ($U_{cp}$) above the array's open-circuit
voltage at the coldest expected temperature with margin, for example
$U_{cp} \geq 1000\,\si{\volt}$ DC for a string whose cold open-circuit
voltage reaches the high $600$s to low $800$s of volts. Choose a Type 2
device (indirect surges) with nominal discharge current $I_n$ of
$20\,\si{\kilo\ampere}$ ($8/20\,\si{\micro\second}$ waveform) and a
maximum discharge current $I_{max}$ of $40\,\si{\kilo\ampere}$. The
voltage protection level $U_p$ should sit comfortably below the
inverter's DC-input surge withstand, typically $U_p \leq 4\,\si{\kilo
\volt}$, and lower if the inverter datasheet demands it. Response time
for a metal-oxide-varistor device is under $25\,\si{\nano\second}$. The
device must be bonded with a short, low-inductance lead to the same
grounding electrode as the array frame, because lead inductance adds
$L\,\dd I/\dd t$ to $U_p$ and can dominate at the steep surge fronts.
The governing standard is IEC 61643-31 for photovoltaic SPDs (with the
installation requirements in IEC 60364-7-712).
\end{exercise}

\subsection*{Diagnosis and reverse engineering}\pathtag{standard}

\begin{exercise}
A magnetic-field reading at the bench shows $300\,\si{\micro\tesla}$
at $20\,\si{\centi\meter}$ from a thick power cable. Identify the
current flowing in the cable, and identify whether the cable carries
DC or AC by what additional measurement would distinguish them.

\paragraph{Solution sketch.} Invert $B = \mu_0 I/(2\pi r)$: $I = 2\pi r
B/\mu_0 = 2\pi(0.2)(300\times 10^{-6})/(4\pi\times 10^{-7}) = 300\,\si{
\ampere}$. The DC-versus-AC question is answered by the probe's
response over time: a DC current gives a steady reading on a Hall probe
and zero reading on a clamp meter or search coil that responds only to
$\dd B/\dd t$; an AC current gives a reading that oscillates at the line
frequency and is detected by an inductive pickup. So the distinguishing
measurement is either a search-coil pickup (responds to AC only) or
watching the Hall probe for a steady versus oscillating value. Note the
clamp around a two-conductor cable carrying balanced go-and-return
current would read near zero; the $300\,\si{\micro\tesla}$ reading
implies a net unbalanced current, such as a single conductor or a fault.
\end{exercise}

\begin{exercise}
A capacitor stamped ``$10\,\si{\micro\farad}, 25\,\si{\volt}$'' is
measured at $9.7\,\si{\micro\farad}$ on an LCR meter at $1\,\si{\kilo
\hertz}$ and $8.4\,\si{\micro\farad}$ at $100\,\si{\kilo\hertz}$.
Identify the most likely cause of the frequency dependence, and the
secondary measurement that would distinguish dielectric absorption
from equivalent series resistance.

\paragraph{Solution sketch.} The apparent capacitance falling with
frequency points to the device's parasitics: a real capacitor is a
series $R$-$L$-$C$ network, and the equivalent series inductance plus
the dielectric's frequency-dependent permittivity both reduce the
measured value as frequency rises. For an electrolytic at $100\,\si{
\kilo\hertz}$ the most likely cause is the combination of equivalent
series resistance (ESR) and inductance shifting the impedance away from
ideal $1/(\omega C)$, together with the polar dielectric's relaxation.
To separate dielectric absorption from ESR, measure the loss tangent or
dissipation factor versus frequency: ESR produces a loss that rises with
frequency as $\omega C \cdot \text{ESR}$, while dielectric relaxation
produces a loss peak at the relaxation frequency. A complementary test
is a dielectric-absorption ("soakage") measurement: charge, briefly
discharge, then watch the voltage recover; the recovery isolates the
slow dielectric mechanism from the resistive ESR.
\end{exercise}

\begin{exercise}
A coaxial cable in a test fixture shows an unexpectedly high VSWR.
Identify three candidate causes (connector mismatch, dielectric
damage, length-and-cut frequency interaction), and the
time-domain-reflectometry signature that distinguishes them.

\paragraph{Solution sketch.} A high voltage standing-wave ratio means
power is being reflected by an impedance discontinuity somewhere along
the line. Time-domain reflectometry (TDR) sends a fast step and watches
the echo; the position of the reflection in time, multiplied by the
propagation velocity, locates the fault. A connector mismatch shows a
sharp localised reflection at the connector position, whose sign tells
whether the discontinuity is high or low impedance. Dielectric damage
(crush, water ingress) shows a reflection partway along the cable at the
damage point, often broadened and with a velocity change beyond it.
A length-and-cut frequency interaction is not a localised fault at all:
the reflection comes from the far-end load, and the VSWR varies
periodically with frequency as the round-trip electrical length passes
through half-wavelengths, with no internal TDR echo. So the
distinguishing signature is the time position and shape of the echo:
at a connector, partway along the dielectric, or only at the far end.
\end{exercise}

\subsection*{Failure analysis}\pathtag{core}

\begin{exercise}
A laboratory data-acquisition system shows $60\,\si{\hertz}$ hum on
every channel even when the input is shorted at the connector.
Identify the failure mode and three measurements that would
distinguish a ground-loop cause from a power-supply-coupling cause.

\paragraph{Solution sketch.} Hum present on every channel with the input
shorted means the interference is entering after the input, through the
shared ground or the power supply rather than through the signal path.
The most common cause is a ground loop: two separate earth paths
enclosing area through which mains-frequency flux induces an EMF.
Three discriminating measurements: (1) lift the chassis earth through an
isolation transformer (safely, briefly) and watch whether the hum drops
sharply, which implicates a ground loop; (2) measure the hum frequency
content, where a pure $60\,\si{\hertz}$ tone suggests inductive
ground-loop pickup while strong harmonics ($120\,\si{\hertz}$ and up)
suggest rectifier ripple from the power supply; (3) probe the supply
rails directly for ripple at $120\,\si{\hertz}$, which appears in
supply-coupling but not in a clean ground loop. The cure for the loop is
a single-point ground; the cure for supply coupling is better filtering
and regulation.
\end{exercise}

\begin{exercise}
A photovoltaic inverter fails after a thunderstorm with no direct
strike to the array. Identify the failure mechanism, the design
defect that allowed it, and the IEC standard that addresses it.

\paragraph{Solution sketch.} The mechanism is induced surge: a nearby
lightning stroke produces a fast-changing magnetic field that links the
loop area formed by the long DC cabling between modules and inverter,
inducing an EMF (Faraday's law) that drives a current exceeding the
inverter's input surge rating, even with no direct strike. The design
defect is the absence of surge protection at the inverter DC input
combined with a large, unmanaged DC-cable loop area and inadequate
bonding of the array frame to the inverter grounding. The cure is a
Type 2 SPD at the inverter DC input, bonded to the same grounding
electrode as the frame, plus routing the DC positive and negative
conductors close together to shrink the loop area. The governing
standards are IEC 61643-31 (photovoltaic SPDs) and the lightning
protection framework IEC 62305; this chapter narrates the mechanism only
and names no specific accident.
\end{exercise}

\begin{exercise}
An MRI suite reports inconsistent image quality correlated with the
time of day. Identify two electromagnetic causes (DC-traction-power
disturbance from a nearby tram line, lift-motor commutation during
business hours) and the measurement that distinguishes them.

\paragraph{Solution sketch.} MRI imaging is acutely sensitive to small
fluctuations in the static field and to stray AC fields in the imaging
band. A nearby tram or light-rail line carries large DC traction
currents whose return paths shift through the ground as trains move,
producing slow, large-amplitude DC field excursions that track the
transit schedule. Lift (elevator) motors commutate when cars run during
business hours, producing AC and broadband disturbances correlated with
occupancy. The distinguishing measurement is a continuous fluxgate (DC-
capable) magnetometer log in the scanner room: the tram signature is a
slow, low-frequency DC drift correlated with the transit timetable,
while the lift signature is intermittent higher-frequency bursts
correlated with car movement. Spectral analysis of the log separates the
near-DC drift from the motor-commutation bursts, and correlating each
against the tram timetable and the lift call log identifies the source.
\end{exercise}

\begin{exercise}
A radio receiver's sensitivity drops by $20\,\si{\decibel}$ when a
nearby LED desk lamp is switched on. Identify the failure mode, the
regulatory limit that should have prevented it, and the
electromagnetic-emissions test that the lamp manufacturer should have
performed.

\paragraph{Solution sketch.} The failure mode is conducted and radiated
EMI from the lamp's switching driver: a cheap LED supply switches at
tens to hundreds of kilohertz with fast edges whose harmonics extend
into the radio bands, raising the receiver's noise floor and cutting its
sensitivity. The regulatory limit that should have prevented it is the
conducted- and radiated-emission limit for lighting equipment, in the
EU the harmonised standard EN 55015 (CISPR 15) under the EMC Directive
2014/30/EU, and in the United States FCC Part 15 for unintentional
radiators. The test the manufacturer should have run is a CISPR 15
conducted-emissions scan on the mains terminals plus a radiated-emission
scan over the relevant frequency range in a calibrated setup (line
impedance stabilisation network for conducted, an antenna at a fixed
distance for radiated). A compliant driver uses input filtering and
edge-rate control to keep harmonics under the limit line. The chapter
narrates the mechanism and the standards; it names no specific product.
\end{exercise}

\subsection*{Judgment}\pathtag{standard}

\begin{exercise}
A junior engineer proposes to route a high-current power cable and a
low-level signal cable through the same conduit ``to save trays.''
Identify the EMC argument against the proposal and the conditions
under which the practice would nevertheless be acceptable.

\paragraph{Solution sketch.} The argument against is crosstalk: the
high-current cable's changing magnetic field couples inductively into
the signal cable's loop, and its changing voltage couples capacitively,
injecting interference proportional to the source's $\dd I/\dd t$ and
$\dd V/\dd t$ and to the run length over which the two share the
conduit. A low-level signal is the worst victim because the coupled
noise is large relative to the wanted signal. The practice becomes
acceptable when the coupling is made negligible: the signal is balanced
(differential) and shielded so common-mode pickup is rejected; the
signal is twisted pair so the loop area is tiny; the power cable carries
DC or low, slowly varying current with small $\dd I/\dd t$; physical
separation and a grounded barrier inside the conduit are maintained; and
the run is short. Many cabling codes set a minimum separation between
power and signal that, if met, permits shared routing. The honest answer
to the junior engineer is conditional, not a flat refusal.
\end{exercise}

\begin{exercise}
A consulting client asks whether they should use a single large
lightning rod at the apex of a $30\,\si{\meter}$ tall industrial
building or a grid of smaller rods distributed across the roof.
Identify the IEC 62305 rolling-sphere model and explain which
approach the model favours.

\paragraph{Solution sketch.} The rolling-sphere method of IEC 62305 models
the lightning attachment point geometrically: a sphere of a defined
radius (set by the protection level, for example $20$ to $60\,\si{
\meter}$) is rolled over the structure, and any surface the sphere
touches is a possible strike point, while any region the sphere cannot
reach is protected. A single rod at the apex of a $30\,\si{\meter}$
building leaves much of the roof exposed, because the sphere rolling
down from the rod still touches the roof edges and any equipment set
back from the apex. A distributed grid of rods, sized and spaced so the
sphere cannot dip to the roof surface between them, protects the whole
roof. The model favours the distributed grid for a large flat or
equipment-laden roof; the single tall rod protects only a cone-like zone
beneath it and is adequate only for a small footprint. The judgment is
to apply the rolling-sphere geometry to the actual roof plan rather than
to argue from the height of a single rod.
\end{exercise}

\begin{exercise}
A regulator is considering raising EMC-emission limits for consumer
electronics to support 5G mobile-radio rollout. Identify the trade-
off between consumer-product cost and radio-service protection, and
the working compromise that has been adopted internationally.

\paragraph{Solution sketch.} Raising emission limits lets manufacturers
spend less on filtering and shielding, lowering product cost, at the
price of a higher ambient noise floor that degrades sensitive radio
services such as mobile reception. Tightening limits protects radio at
the cost of more expensive products. The working compromise adopted
internationally is frequency-dependent, banded limit lines: emission
limits are set per frequency band, stricter in the bands allocated to
protected radio services and looser elsewhere, with separate conducted
and radiated limits and defined measurement methods (the CISPR family
of standards). The regulator does not pick a single number; it shapes
the limit curve to protect the spectrum where protection matters while
relaxing it where no sensitive service operates. The judgment is to
frame the question as a per-band limit curve, not a single global
threshold.
\end{exercise}

\begin{exercise}
An industrial facility experiences periodic motor-drive failures
attributed to ``power quality.'' A consultant proposes installing
$\$50\,000$ of harmonic filters; another proposes $\$5\,000$ of
isolation transformers. Identify what measurements would settle the
choice and the conditions under which each remedy is appropriate.

\paragraph{Solution sketch.} The measurement that settles it is a power-
quality survey at the drive's supply: log the voltage and current
waveforms and compute the harmonic spectrum (total harmonic distortion
and the individual harmonic magnitudes), the flicker, and any transient
events, over a representative window that captures the failure pattern.
If the dominant problem is harmonic current distortion (large fifth,
seventh, eleventh harmonics from nonlinear loads), harmonic filters are
the appropriate remedy and the $\$50{,}000$ spend may be justified. If
the problem is common-mode noise, ground-loop coupling, or a need to
break a galvanic path, an isolation transformer at $\$5{,}000$ solves it
and the expensive filters would not. If the failures correlate with
voltage sags or transients from other equipment, neither remedy fits and
the cure is a line conditioner or upstream fault correction. The honest
engineering answer is to measure before specifying: the spectrum and the
event log tell which mechanism is operating, and each remedy treats only
its own mechanism.
\end{exercise}
